% Chapter 4, Section 1 _Linear Algebra_ Jim Hefferon
%  http://joshua.smcvt.edu/linearalgebra
%  2001-Jun-12
\chapter{Determinan} 
Dalam bab pertama kita  
menyoroti kasus khusus sistem linear
yang memiliki jumlah persamaan sama dengan jumlah variabel tak diketahui, 
yaitu yang berbentuk \( T\vec{x}=\vec{b} \), dengan $T$ suatu matriks persegi.
Kita mencatat bahwa hanya ada dua jenis $T$.
Jika $T$ berkaitan dengan solusi tunggal
untuk suatu $\vec{b}$, 
seperti pada sistem homogen $T\vec{x}=\zero$, maka
$T$ berkaitan dengan solusi tunggal untuk setiap $\vec{b}$ tersebut.
Matriks seperti itu kita sebut nonsingular.
Jenis $T$ lainnya, yakni yang setiap sistem linearnya---dengan matriks tersebut sebagai
matriks koefisien---tidak memiliki solusi atau memiliki tak hingga banyak solusi,
kita sebut singular.  

Dalam pembahasan selanjutnya, pembedaan ini terus menjadi tema utama.
Sebagai contoh, kini kita mengetahui bahwa matriks 
\( \nbyn{n} \) \( T \) bersifat nonsingular jika dan hanya jika
semua pernyataan berikut berlaku:
%<*EquivalentOfNonsingular>
\begin{itemize}
   \item setiap sistem \( T\vec{x}=\vec{b} \) memiliki solusi 
          dan solusi itu tunggal;
   \item reduksi Gauss--Jordan terhadap $T$ menghasilkan matriks identitas;
   \item baris-baris $T$ membentuk himpunan bebas linear;
   \item kolom-kolom \( T \) membentuk himpunan bebas linear, 
         yakni suatu basis bagi \( \Re^n \);
   \item setiap pemetaan yang direpresentasikan oleh \( T \) merupakan isomorfisme;
   \item terdapat matriks invers \( T^{-1} \).
\end{itemize}
%</EquivalentOfNonsingular>
Karena itu, ketika mengamati sebuah matriks persegi, salah satu hal pertama
yang kita tanyakan adalah apakah matriks tersebut  
nonsingular.

Bab ini mengembangkan suatu rumus untuk menentukan apakah $T$ nonsingular.
Lebih tepatnya, kita akan mengembangkan suatu rumus 
untuk matriks $\nbyn{1}$, suatu rumus untuk matriks $\nbyn{2}$, dan seterusnya.
Rumus-rumus ini saling berkaitan secara alami; dengan kata lain, 
kita akan mengembangkan suatu keluarga
rumus, yaitu suatu skema yang menjelaskan rumus untuk setiap ukuran.

Karena pembahasan kita akan dibatasi pada matriks persegi, dalam bab ini
kita akan sering cukup mengatakan `matriks' sebagai pengganti `matriks persegi'.





\section{Def{}inisi}
%<*DeterminantIntro>
Menentukan kenonsingularan
sangat mudah untuk matriks \( \nbyn{1} \).
\begin{equation*}
  \begin{mat}
       a
  \end{mat} 
  \quad\text{nonsingular jika dan hanya jika}\quad
  a \neq 0
\end{equation*} 
Korolari~Tiga.IV.\ref{cor:TwoByTwoInv} memberikan rumus untuk ukuran $\nbyn{2}$. 
\begin{equation*}
     \begin{mat}
           a  &b  \\
           c  &d
      \end{mat}  
   \quad\text{nonsingular jika dan hanya jika}\quad
   ad-bc \neq 0
\end{equation*}
Kita dapat memperoleh rumus untuk ukuran $\nbyn{3}$ seperti memperoleh rumus sebelumnya,  
meskipun perhitungannya rumit
% \typeout{DET1 ThreeByThreeDetForm cleveref expmt!}
% (see \cref{exer:ThreeByThreeDetForm}).
(see \nearbyexercise{exer:ThreeByThreeDetForm}).
\begin{equation*}
     \begin{mat}
           a  &b  &c  \\
           d  &e  &f  \\
           g  &h  &i
     \end{mat}
  \quad\text{nonsingular jika dan hanya jika}\quad
    aei+bfg+cdh-hfa-idb-gec \neq 0
\end{equation*}
Dengan mengingat kasus-kasus ini, kita mengajukan suatu keluarga 
rumus: $a$, $ad-bc$, dan seterusnya. 
Untuk setiap $n$, rumus tersebut mendefinisikan suatu 
\definend{determinan}\index{determinan}\index{matriks!determinan} 
fungsi
$\map{\det_{\nbyn{n}}}{\matspace_{\nbyn{n}}}{\Re}$ 
sedemikian sehingga matriks $\nbyn{n}$ $T$ nonsingular jika dan
hanya jika $\det_{\nbyn{n}}(T)\neq 0$.
%</DeterminantIntro>
(Biasanya kita menghilangkan subskrip \( \nbyn{n} \) karena 
ukuran \( T \) menunjukkan fungsi determinan mana yang dimaksud.)






\subsectionoptional{Eksplorasi}
\textit{Subbagian ini merupakan motivasi dan pengembangan opsional
untuk definisi umum.
Definisi tersebut terdapat dalam subbagian berikutnya.}

Di atas, pada setiap kasus matriks bersifat nonsingular jika dan hanya 
jika suatu rumus bernilai tak nol.
Namun, ketiga rumus tersebut tidak 
memperlihatkan pola yang jelas. 
Kita mungkin memperhatikan bahwa suku \(\nbyn{1}\)
\( a \) memiliki satu huruf, suku-suku \(\nbyn{2}\)
\(ad\) dan \(bc\) memiliki dua huruf, dan setiap suku \(\nbyn{3}\)
memiliki tiga huruf.
Kita bahkan mungkin memperhatikan bahwa di dalam suku-suku tersebut
terdapat satu huruf dari setiap baris dan setiap kolom matriks; misalnya, 
dalam suku \(cdh\), masing-masing satu huruf 
berasal dari setiap baris dan setiap kolom.
\begin{equation*}
   \begin{mat}
          &    &c \\
      d           \\
          &h
   \end{mat} 
\end{equation*}
Namun, pengamatan ini mungkin lebih membingungkan daripada
mencerahkan.
Sebagai contoh, kita mungkin bertanya-tanya mengapa 
sebagian suku dijumlahkan, sedangkan sebagian lainnya dikurangkan.

Strategi yang baik untuk menyelesaikan masalah ialah
menyelidiki sifat apa saja yang harus dimiliki solusinya, lalu
mencari sesuatu yang memiliki sifat-sifat tersebut.
Karena itu, kita akan mulai dengan menanyakan sifat apa yang kita inginkan
dimiliki oleh rumus-rumus determinan.

Pada tahap ini, cara utama kita
untuk menentukan apakah suatu matriks singular atau tidak
ialah melakukan reduksi 
Gauss, lalu memeriksa apakah 
diagonal matriks bentuk eselon memuat nol,
yakni apakah hasil kali entri-entri diagonalnya bernilai~nol.
Jadi, apa pun rumus determinan yang kita temukan, kita dapat menduga bahwa bukti 
kebenarannya mungkin melibatkan penerapan Metode Gauss pada matriks tersebut
untuk menunjukkan bahwa pada akhirnya hasil kali entri-entri diagonal bernilai nol jika dan hanya jika
rumus kita memberikan nol. 

Hal ini menyarankan suatu rencana:~kita akan mencari keluarga rumus
determinan yang 
tidak berubah oleh operasi baris dan sedemikian sehingga determinan suatu
matriks bentuk eselon merupakan hasil kali entri-entri diagonalnya.
% Under this plan, a proof that the functions determine singularity would go, 
% ``Where $T\rightarrow\cdots\rightarrow\hat{T}$ is the Gaussian
% reduction, the determinant of $T$ equals the
% determinant of $\hat{T}$ (because the determinant is unchanged by row
% operations), which is the product down the diagonal, which is
% zero if and only if the matrix is singular''.
Dalam sisa subbagian ini, kita akan menguji rencana tersebut terhadap 
rumus $\nbyn{2}$ dan $\nbyn{3}$.
Pada akhirnya kita perlu mengubah bagian ``tidak berubah oleh operasi baris'', 
tetapi perubahannya tidak banyak.

Pertama-tama kita memeriksa apakah   
rumus $\nbyn{2}$ dan $\nbyn{3}$ tidak berubah oleh 
operasi kombinasi baris:~jika 
\begin{equation*}
   T \grstep{k\rho_i+\rho_j} \hat{T}
\end{equation*}
apakah kemudian \( \det(\hat{T})=\det(T) \)? 
Pemeriksaan determinan $\nbyn{2}$ setelah operasi $k\rho_1+\rho_2$
\begin{equation*}
   \det(
       \begin{mat}
         a     &b       \\
         ka+c  &kb+d    \\
      \end{mat}
   )
   = a(kb+d)-(ka+c)b = ad-bc
\end{equation*}
menunjukkan bahwa nilainya memang tidak berubah, dan
kombinasi $\nbyn{2}$ lainnya, $k\rho_2+\rho_1$, memberikan hasil yang sama.
Demikian pula, 
kombinasi $\nbyn{3}$ $k\rho_3+\rho_2$ tidak mengubah determinan
  \begin{align*}
    \det(
      \begin{mat}
         a    &b    &c    \\
         kg+d &kh+e &ki+f \\
         g    &h    &i
      \end{mat}
    )
    &=\begin{array}[t]{@{}l@{}}
         a(kh+e)i+b(ki+f)g+c(kg+d)h \\
         \ \hbox{}-h(ki+f)a-i(kg+d)b-g(kh+e)c  
    \end{array}                                 \\
     &=aei + bfg + cdh - hfa - idb - gec
  \end{align*}
sebagaimana juga operasi kombinasi baris $\nbyn{3}$ lainnya.

Jadi, rencana tersebut tampaknya menjanjikan.
Tentu saja, seandainya kita telah memperoleh 
rumus determinan $\nbyn{4}$ dan mengujinya, mungkin saja kita menemukan
bahwa rumus itu berubah oleh kombinasi baris.
Ini merupakan sebuah eksplorasi dan kita belum memiliki semua faktanya.
Meskipun demikian, sejauh ini hasilnya baik.

Selanjutnya kita membandingkan \( \det(\hat{T}) \) dengan 
\( \det(T) \) untuk pertukaran baris.
% \begin{equation*}
%    T \grstep{ {\rho}_i \leftrightarrow {\rho}_j } \hat{T}
% \end{equation*}
Di sini kita menemui kendala:
pertukaran baris \(\nbyn{2}\) $\rho_1\leftrightarrow\rho_2$
tidak menghasilkan \( ad-bc \).
  \begin{equation*}
     \det(
       \begin{mat}
         c  &d \\
         a  &b
       \end{mat}
     )
     = bc - ad
  \end{equation*}
Pertukaran $\rho_1\leftrightarrow\rho_3$ di dalam matriks \(\nbyn{3}\) ini
\begin{equation*}
   \det(
     \begin{mat}
        g  &h  &i \\
        d  &e  &f \\
        a  &b  &c
     \end{mat}
   )
   = gec + hfa + idb - bfg - cdh - aei
\end{equation*}
juga tidak memberikan determinan yang sama seperti sebelum pertukaran karena sekali lagi 
terjadi perubahan tanda.
Mencoba pertukaran \(\nbyn{3}\) lain, $\rho_1\leftrightarrow\rho_2$, 
\begin{equation*}
   \det(
     \begin{mat}
        d  &e  &f \\
        a  &b  &c \\
        g  &h  &i
     \end{mat}
   )
   = dbi + ecg + fah - hcd - iae - gbf
\end{equation*}
juga menghasilkan perubahan tanda.

Jadi, dalam percobaan ini pertukaran baris tampaknya  
mengubah tanda determinan.
Hal ini tidak sepenuhnya merusak rencana kita.
Kita berharap dapat menentukan kenonsingularan hanya dengan mempertimbangkan
apakah rumus tersebut memberikan nol, bukan dengan mempertimbangkan tandanya.
Karena itu, alih-alih mengharapkan rumus determinan sama sekali
tidak berubah oleh operasi baris, kita mengubah rencana agar rumus tersebut
berganti tanda ketika terjadi pertukaran.

Sebagai langkah terakhir, tentu saja
kita membandingkan \( \det(\hat{T}) \) dengan \( \det(T) \) 
untuk operasi
% \begin{equation*}
%    T \grstep{ k{\rho}_i } \hat{T}
% \end{equation*}
perkalian suatu baris dengan skalar. 
Perhitungan ini
\begin{equation*}
  \det(
    \begin{mat}
      a   &b   \\
      kc  &kd
    \end{mat}
  )
  = a(kd) - (kc)b
  =k\cdot (ad-bc)
\end{equation*}
menghasilkan seluruh determinan dikalikan dengan~$k$,
dan kasus $\nbyn{2}$ lainnya memberikan hasil yang sama.
Kasus \(\nbyn{3}\) ini berakhir dengan cara yang sama
\begin{align*}
  \det(
  \begin{mat}
    a    &b    &c   \\
    d    &e    &f   \\
    kg   &kh   &ki
  \end{mat}
  )
   &= \begin{array}[t]{@{}l@{}}
         ae(ki) + bf(kg) + cd(kh)                \\
         \>- (kh)fa - (ki)db - (kg)ec  
      \end{array}                                      \\
   &= k\cdot(aei + bfg + cdh - hfa - idb - gec)
\end{align*}
sebagaimana juga dua kasus $\nbyn{3}$ lainnya.
Perhitungan-perhitungan ini membuat kita menduga bahwa mengalikan suatu baris dengan~$k$
akan mengalikan determinannya dengan~$k$.
Seperti sebelumnya, hal ini mengubah rencana kita, tetapi tidak merusaknya.
Kita hanya mensyaratkan agar
sifat bernilai nol dari rumus determinan tidak berubah, tanpa berfokus pada
tanda atau besarnya.

Jadi, dalam eksplorasi ini rencana kita 
mengalami beberapa perubahan yang tidak mendasar dan sekarang berbunyi:~kita akan mencari 
fungsi determinan $\nbyn{n}$
yang tetap tidak berubah
di bawah operasi kombinasi baris, yang berganti tanda ketika
dua baris dipertukarkan, yang ikut terskala ketika suatu baris diskalakan,
dan sedemikian sehingga determinan suatu matriks bentuk eselon merupakan
hasil kali entri-entri diagonalnya.
Dalam dua subbagian berikutnya kita akan melihat bahwa untuk setiap~$n$
terdapat tepat satu fungsi semacam itu.

Terakhir, untuk subbagian berikutnya, perhatikan bahwa mengeluarkan faktor skalar merupakan 
operasi per baris:
di sini
\begin{equation*}
    \det(
      \begin{mat}[r]
        3  &3  &9  \\
        2  &1  &1  \\
        5  &11 &-5
     \end{mat}
    )
    =3 \cdot \det(
               \begin{mat}[r]
                 1  &1  &3  \\
                 2  &1  &1  \\
                 5  &11 &-5
               \end{mat}
             )                                  
\end{equation*}
faktor $3$ hanya dikeluarkan dari baris teratas, sehingga baris-baris lainnya tidak berubah.
Akibatnya, dalam definisi determinan kita akan
menuliskannya sebagai fungsi dari baris-baris
\( \det (\vec{\rho}_1,\vec{\rho}_2,\dots\vec{\rho}_n) \), bukan sebagai 
\( \det(T) \) atau sebagai fungsi entri-entri
\( \det(t_{1,1},\dots,t_{n,n}) \).

\begin{exercises}
  \recommended \item 
     Hitung determinan masing-masing matriks berikut.
     \begin{exparts*}
       \partsitem \(
            \begin{mat}[r]
                3    &1   \\
               -1    &1
            \end{mat}    \)
       \partsitem \(
            \begin{mat}[r]
                2    &0   &1  \\
                3    &1   &1 \\
               -1    &0   &1
            \end{mat}    \)
       \partsitem \(
            \begin{mat}[r]
                4    &0   &1  \\
                0    &0   &1 \\
                1    &3   &-1
            \end{mat}    \)
     \end{exparts*}
     \begin{answer}
       \begin{exparts*}
         \partsitem \( 4 \)
         \partsitem \( 3 \)
         \partsitem \( -12 \)
       \end{exparts*}  
     \end{answer}
  \item 
     Hitung determinan masing-masing matriks berikut.
     \begin{exparts*}
        \partsitem \( \begin{mat}[r]
                    2  &0  \\
                   -1  &3
                 \end{mat} \)
        \partsitem \( \begin{mat}[r]
                    2  &1  &1  \\
                    0  &5  &-2 \\
                    1  &-3 &4
                 \end{mat} \)
        \partsitem \( \begin{mat}[r]
                    2  &3  &4  \\
                    5  &6  &7  \\
                    8  &9  &1
                 \end{mat} \)
     \end{exparts*}
    \begin{answer}
      \begin{exparts*}
        \partsitem \( 6 \)
        \partsitem \( 21 \)
        \partsitem \( 27 \)
      \end{exparts*}  
    \end{answer}
  \recommended \item  
    Verifikasikan bahwa determinan matriks segitiga atas
    $\nbyn{3}$ merupakan hasil kali entri-entri diagonalnya. 
    \begin{equation*}
       \det(
       \begin{mat}
           a    &b   &c    \\
           0    &e   &f    \\
           0    &0   &i
       \end{mat}
       )
       =aei
    \end{equation*}
    Apakah matriks segitiga bawah berlaku dengan cara yang sama?
    \begin{answer}
      Untuk yang pertama, terapkan rumus dalam bagian ini, perhatikan bahwa setiap
      suku yang memuat \( d \), \( g \), atau \( h \) bernilai nol, lalu sederhanakan.
      Matriks segitiga bawah berlaku dengan cara yang sama.  
    \end{answer}
  \recommended \item 
     Gunakan determinan untuk menentukan apakah masing-masing matriks singular atau
     nonsingular.
     \begin{exparts*}
       \partsitem \(
            \begin{mat}[r]
                2    &1   \\
                3    &1
            \end{mat}    \)
       \partsitem \(
            \begin{mat}[r]
                0    &1   \\
                1    &-1
            \end{mat}    \)
       \partsitem \(
            \begin{mat}[r]
                4    &2   \\
                2    &1
            \end{mat}    \)
     \end{exparts*}
     \begin{answer}
       \begin{exparts}
         \partsitem Nonsingular, determinannya adalah \( -1 \).
         \partsitem Nonsingular, determinannya adalah \( -1 \).
         \partsitem Singular, determinannya adalah \( 0 \).
       \end{exparts}  
     \end{answer}
  \item 
     Singular atau nonsingular?
     Gunakan determinan untuk menentukannya.
     \begin{exparts*}
       \partsitem \(
            \begin{mat}[r]
                2    &1   &1  \\
                3    &2   &2 \\
                0    &1   &4
            \end{mat}    \)
       \partsitem \(
            \begin{mat}[r]
                1    &0   &1  \\
                2    &1   &1 \\
                4    &1   &3
            \end{mat}    \)
       \partsitem \(
            \begin{mat}[r]
                2    &1   &0  \\
                3    &-2  &0 \\
                1    &0   &0
            \end{mat}    \)
     \end{exparts*}
     \begin{answer}
      \begin{exparts}
         \partsitem Nonsingular, determinannya adalah \( 3 \).
         \partsitem Singular, determinannya adalah \( 0 \).
         \partsitem Singular, determinannya adalah \( 0 \).
       \end{exparts}  
     \end{answer}
  \recommended \item
    Setiap pasangan matriks berbeda oleh satu operasi baris.
    Gunakan operasi tersebut untuk
    membandingkan \( \det(A) \) dengan \( \det(B) \).
     \begin{exparts}
        \partsitem \( A=\begin{mat}[r]
                      1  &2  \\
                      2  &3
                   \end{mat} \), 
               \(  B=\begin{mat}[r]
                      1  &2  \\
                      0  &-1
                   \end{mat}  \)
        \partsitem \( A=\begin{mat}[r]
                      3  &1  &0  \\
                      0  &0  &1  \\
                      0  &1  &2
                   \end{mat} \),
              \(   B=\begin{mat}[r]
                      3  &1  &0  \\
                      0  &1  &2  \\
                      0  &0  &1
                   \end{mat}  \)
        \partsitem \( A=\begin{mat}[r]
                      1  &-1 &3  \\
                      2  &2  &-6 \\
                      1  &0  &4
                   \end{mat} \),
              \(   B=\begin{mat}[r]
                      1  &-1 &3  \\
                      1  &1  &-3 \\
                      1  &0  &4
                   \end{mat}  \)
     \end{exparts}
     \begin{answer}
       \begin{exparts}
         \partsitem \( \det(B)=\det(A) \) melalui \( -2\rho_1+\rho_2 \)
         \partsitem \( \det(B)=-\det(A) \) melalui 
             \( \rho_2\leftrightarrow\rho_3 \)
         \partsitem \( \det(B)=(1/2)\cdot \det(A) \) melalui \( (1/2)\rho_2 \)
       \end{exparts}   
     \end{answer}
  \recommended \item Tentukan determinan matriks $\nbyn{4}$ berikut 
    dengan mengikuti rencana ini:~lakukan Metode Gauss dan 
    perhatikan bahwa determinan tetap tidak berubah
    pada kombinasi baris, berganti tanda pada
    pertukaran baris, ikut terskala ketika suatu baris diskalakan,
    dan bahwa determinan matriks bentuk eselon merupakan
    hasil kali entri-entri diagonalnya.
    \begin{equation*}
      \begin{mat}
        1 &2 &0  &2 \\
        2 &4 &1  &0 \\
        0 &0 &-1 &3 \\
        3 &-1 &1 &4
      \end{mat}
    \end{equation*}
    \begin{answer}
      Metode Gauss memberikan hasil berikut.
      % sage: m = matrix(QQ, [[1,2,0,2], [2,4,1,0], [0,0,-1,3], [3,-1,1,4]])
      % sage: gauss_method(m)
      % [ 1  2  0  2]
      % [ 2  4  1  0]
      % [ 0  0 -1  3]
      % [ 3 -1  1  4]
      %  take -2 times row 1 plus row 2
      %  take -3 times row 1 plus row 4
      % [ 1  2  0  2]
      % [ 0  0  1 -4]
      % [ 0  0 -1  3]
      % [ 0 -7  1 -2]
      %  swap row 2 with row 4
      % [ 1  2  0  2]
      % [ 0 -7  1 -2]
      % [ 0  0 -1  3]
      % [ 0  0  1 -4]
      %  take 1 times row 3 plus row 4
      % [ 1  2  0  2]
      % [ 0 -7  1 -2]
      % [ 0  0 -1  3]
      % [ 0  0  0 -1]
      \begin{equation*}
        \begin{mat}
          1 &2 &0  &2 \\
          2 &4 &1  &0 \\
          0 &0 &-1 &3 \\
          3 &-1 &1 &4
        \end{mat}
        \grstep[-3\rho_1+\rho_4]{-2\rho_1+\rho_2}
        \grstep{\rho_2\leftrightarrow\rho_4}
        \grstep{\rho_3+\rho_4}      
        \begin{mat}
          1 &2 &0  &2 \\
          0 &-7 &1  &-2 \\
          0 &0 &-1 &3 \\
          0 &0 &0 &-1
        \end{mat}
      \end{equation*}
      Hasil kali entri-entri diagonal matriks bentuk eselon adalah 
       $1\cdot(-7)\cdot(-1)\cdot(-1)=-7$.
      Selama Metode Gauss tidak ada baris yang diskalakan, tetapi terjadi 
      satu pertukaran baris. Karena itu, untuk memperoleh determinan kita mengganti tandanya,
      sehingga diperoleh $+7$.
    \end{answer}
  \item 
     Tunjukkan pernyataan berikut.
      \begin{equation*}
         \det(
         \begin{mat}
             1    &1   &1    \\
             a    &b   &c    \\
             a^2  &b^2 &c^2
         \end{mat}
         )
         =(b-a)(c-a)(c-b)
      \end{equation*}
     \begin{answer}
       Dengan menggunakan rumus determinan matriks $\nbyn{3}$,
       kita uraikan ruas kiri
       \begin{equation*}
         1\cdot b\cdot c^2+1\cdot c\cdot a^2+1\cdot a\cdot b^2
          -b^2\cdot c\cdot 1 -c^2\cdot a\cdot 1-a^2\cdot b\cdot 1
       \end{equation*}
       dan dengan menggunakan sifat distributif kita uraikan ruas kanan.
       \begin{equation*}
         (bc-ba-ac+a^2)\cdot(c-b)
         =c^2b-b^2c-bac+b^2a-ac^2+acb+a^2c-a^2b
       \end{equation*}
       Sekarang kita tinggal memeriksa bahwa keduanya sama.
       (\textit{Catatan}.
       Ini merupakan kasus \( \nbyn{3} \) dari
       \definend{determinan Vandermonde}\index{Vandermonde!determinan}%
       \index{determinan!Vandermonde} yang muncul dalam berbagai penerapan).
     \end{answer}
   \recommended \item 
      Bilangan real \( x \) mana yang membuat matriks ini singular?
      \begin{equation*}
         \begin{mat}
            12-x  &4  \\
            8    &8-x
         \end{mat}
      \end{equation*}
      \begin{answer}
         Persamaan
         \begin{equation*}
           0=
           \det(
             \begin{mat}
                12-x  &4  \\
                8    &8-x
             \end{mat}
           )
           =64-20x+x^2
           =(x-16)(x-4) 
       \end{equation*}
       memiliki akar \( x=16 \) dan \( x=4 \).  
     \end{answer}
  \item \label{exer:ThreeByThreeDetForm} 
    Lakukan reduksi Gauss untuk memeriksa
    rumus matriks $\nbyn{3}$ yang dinyatakan dalam pendahuluan
    bagian ini.
    \begin{center}
      \( \begin{mat}
               a  &b  &c  \\
               d  &e  &f  \\
               g  &h  &i
         \end{mat} \)
      nonsingular jika dan hanya jika
      \( aei+bfg+cdh-hfa-idb-gec \neq 0 \)
    \end{center}
    \begin{answer}
      Pertama-tama kita mereduksi matriks tersebut ke bentuk eselon.
      Sebagai permulaan, anggap bahwa \( a\neq 0 \) dan \( ae-bd\neq 0 \).
      \begin{multline*}
        \grstep{(1/a)\rho_1}
        \begin{mat}
           1   &b/a   &c/a   \\
           d   &e     &f     \\
           g   &h     &i
         \end{mat}                                           
        \grstep[-g\rho_1+\rho_3]{-d\rho_1+\rho_2}
        \begin{mat}
           1   &b/a           &c/a           \\
           0   &(ae-bd)/a     &(af-cd)/a     \\
           0   &(ah-bg)/a     &(ai-cg)/a
         \end{mat}                                            \\
        \grstep{(a/(ae-bd))\rho_2}
        \begin{mat}
           1   &b/a           &c/a             \\
           0   &1             &(af-cd)/(ae-bd) \\
           0   &(ah-bg)/a     &(ai-cg)/a
         \end{mat}
      \end{multline*}
      Langkah berikut menyelesaikan perhitungannya.
      \begin{equation*}
        \grstep{((ah-bg)/a)\rho_2+\rho_3}
        \begin{mat}
           1   &b/a    &c/a             \\
           0   &1      &(af-cd)/(ae-bd)      \\
           0   &0      &(aei+bgf+cdh-hfa-idb-gec)/(ae-bd)
         \end{mat}
      \end{equation*}
      Sekarang, dengan menganggap bahwa $a\neq 0$ dan \( ae-bd\neq 0 \), 
      matriks semula bersifat nonsingular
      jika dan hanya jika entri \( 3,3 \) di atas bernilai tak nol.
      Dengan kata lain, di bawah anggapan tersebut, matriks semula
      nonsingular jika dan hanya jika $aei+bgf+cdh-hfa-idb-gec\neq 0$,
      sebagaimana yang dikehendaki.

      Kita menutup pembahasan dengan menelusuri apa yang terjadi jika anggapan yang
      digunakan demi kemudahan dalam paragraf sebelumnya tidak berlaku.
      Pertama, jika \( a\neq 0 \), tetapi \( ae-bd=0 \), kita dapat melakukan pertukaran
      \begin{equation*}
        \begin{mat}
           1   &b/a           &c/a           \\
           0   &0             &(af-cd)/a     \\
           0   &(ah-bg)/a     &(ai-cg)/a
         \end{mat}                                  
        \grstep{\rho_2\leftrightarrow\rho_3}
        \begin{mat}
           1   &b/a           &c/a           \\
           0   &(ah-bg)/a     &(ai-cg)/a     \\
           0   &0             &(af-cd)/a
         \end{mat}
      \end{equation*}
      dan menyimpulkan bahwa matriks tersebut nonsingular jika dan hanya jika
      \( ah-bg\neq 0 \) dan \( af-cd\neq 0 \).
      Kondisi `\( ah-bg\neq 0 \) dan \( af-cd\neq 0 \)' ekuivalen dengan
      kondisi `\( (ah-bg)(af-cd)\neq 0 \)'.
      Dengan mengalikan dan menggunakan anggapan kasus $ae-bd=0$
      untuk menyubstitusikan $ae$ menggantikan $bd$, kita memperoleh hasil berikut.
      \begin{multline*}
         0\neq ahaf-ahcd-bgaf+bgcd
          =ahaf-ahcd-bgaf+aegc            \\
          =a(haf-hcd-bgf+egc)
      \end{multline*}
      Karena \( a\neq 0 \), matriks tersebut
      nonsingular jika dan hanya jika \( haf-hcd-bgf+egc\neq 0 \).
      Jadi, dalam kasus \( a\neq 0 \) dan \( ae-bd=0 \) ini, 
      matriks tersebut nonsingular ketika
      \( haf-hcd-bgf+egc-i(ae-bd)\neq 0 \).

      Kasus-kasus sisanya bersifat rutin.
      Kerjakan kasus \( a=0 \), tetapi \( d\neq 0 \), serta kasus \( a=0 \) dan \( d=0 \),
      tetapi \( g\neq 0 \), dengan terlebih dahulu mempertukarkan baris, lalu melanjutkan seperti
      di atas.
      Kasus \( a=0 \), \( d=0 \), dan \( g=0 \) mudah\Dash matriks itu
      singular karena kolom-kolomnya membentuk himpunan bergantung linear, dan
      determinannya bernilai nol.  
    \end{answer}
  \item 
     Tunjukkan bahwa persamaan garis dalam \( \Re^2 \) yang melalui \( (x_1,y_1) \)
     dan \( (x_2,y_2) \) diberikan oleh determinan berikut.
     \begin{equation*}
        \det(
        \begin{mat}
           x   &y   &1  \\
           x_1 &y_1 &1  \\
           x_2 &y_2 &1
        \end{mat})=0 \qquad x_1\neq x_2
     \end{equation*}
     \begin{answer}
       Menghitung determinan dan melakukan sedikit manipulasi aljabar memberikan hasil berikut.
       \begin{align*}
          0
          &=y_1x+x_2y+x_1y_2-y_2x-x_1y-x_2y_1     \\
          (x_2-x_1)\cdot y
          &=(y_2-y_1)\cdot x+x_2y_1-x_1y_2              \\
          y
          &=\frac{y_2-y_1}{x_2-x_1}\cdot x+\frac{x_2y_1-x_1y_2}{x_2-x_1}
       \end{align*}
       Perhatikan bahwa ini merupakan persamaan suatu garis (khususnya,
       persamaan ini memuat bentuk kemiringan yang sudah dikenal), 
       dan perhatikan bahwa \( (x_1,y_1) \) serta \( (x_2,y_2) \) memenuhinya. 
     \end{answer}
  \item
    Banyak orang telah mempelajari
    jembatan keledai berikut untuk determinan matriks \( \nbyn{3} \):~salin dua kolom pertama ke sebelah kanan matriks, 
    lalu ambil hasil kali sepanjang diagonal turun ke kanan dan jumlahkan,
    lalu ambil hasil kali sepanjang diagonal turun ke kiri dan kurangkan.
    Dengan kata lain, tuliskan terlebih dahulu
    \begin{equation*}
        \begin{pmat}{ccc|cc}
          h_{1,1} &h_{1,2} &h_{1,3} &h_{1,1} &h_{1,2} \\
          h_{2,1} &h_{2,2} &h_{2,3} &h_{2,1} &h_{2,2} \\
          h_{3,1} &h_{3,2} &h_{3,3} &h_{3,1} &h_{3,2}
        \end{pmat}
     \end{equation*}
     lalu hitung bentuk berikut.
     \begin{equation*}
      \begin{array}{l}
      h_{1,1}h_{2,2}h_{3,3}+h_{1,2}h_{2,3}h_{3,1}+h_{1,3}h_{2,1}h_{3,2} \\
      \>-h_{3,1}h_{2,2}h_{1,3}-h_{3,2}h_{2,3}h_{1,1}
        -h_{3,3}h_{2,1}h_{1,2}
      \end{array}
    \end{equation*}
    \begin{exparts}
      \partsitem Periksa bahwa hasil ini sesuai dengan rumus yang diberikan dalam 
        pendahuluan bagian ini.
      \partsitem Apakah cara tersebut dapat diperluas ke determinan berukuran lain?
    \end{exparts}
    \begin{answer}
      \begin{exparts}
        \partsitem Perbandingan dengan rumus yang diberikan dalam pendahuluan
          bagian ini mudah dilakukan.
        \partsitem Meskipun cara tersebut berlaku untuk matriks \( \nbyn{2} \),
          \begin{align*}
            \begin{pmat}{cc|c}
              h_{1,1} &h_{1,2} &h_{1,1} \\
              h_{2,1} &h_{2,2} &h_{2,1}
            \end{pmat}
            &=\begin{array}[t]{@{}l@{}}
              h_{1,1}h_{2,2}+h_{1,2}h_{2,1}  \\
              \>-h_{2,1}h_{1,2}-h_{2,2}h_{1,1}
            \end{array}                                     \\
            &=h_{1,1}h_{2,2}-h_{1,2}h_{2,1}
          \end{align*}
          cara itu tidak berlaku untuk matriks \( \nbyn{4} \).
          Sebagai contoh, matriks berikut
          singular karena baris kedua dan ketiganya sama, 
          \begin{equation*}
            \begin{mat}[r]  
              1  &0  &0  &1   \\
              0  &1  &1  &0   \\
              0  &1  &1  &0   \\
             -1  &0  &0  &1  
            \end{mat}
          \end{equation*}
          tetapi mengikuti skema jembatan keledai tersebut tidak memberikan nol.
          \begin{equation*}
            \begin{pmat}{rrrr|rrr}
                       1  &0  &0  &1  &1  &0  &0  \\
                       0  &1  &1  &0  &0  &1  &1  \\
                       0  &1  &1  &0  &0  &1  &1  \\
                      -1  &0  &0  &1  &-1 &0  &0
            \end{pmat}
            =\begin{array}[t]{@{}l@{}}
               1+0+0+0 \\
               \>-(-1)-0-0-0
              \end{array}
           \end{equation*}
      \end{exparts}  
     \end{answer}
  \item 
    \definend{Hasil kali silang}\index{hasil kali silang}\index{vektor!hasil kali silang} 
dari vektor-vektor
    \begin{equation*}
      \vec{x}=\colvec{x_1 \\ x_2 \\ x_3}
      \qquad
      \vec{y}=\colvec{y_1 \\ y_2 \\ y_3}
    \end{equation*}
    adalah vektor yang dihitung sebagai determinan berikut.
    \begin{equation*}
      \vec{x}\times\vec{y}=
      \det(\begin{mat}
        \vec{e}_1  &\vec{e}_2  &\vec{e}_3  \\
        x_1        &x_2        &x_3        \\
        y_1        &y_2        &y_3
      \end{mat})
    \end{equation*}
    Perhatikan bahwa entri-entri baris pertama merupakan vektor, yakni vektor-vektor dari
    basis standar untuk $\Re^3$.
    Tunjukkan bahwa hasil kali silang dua vektor tegak lurus terhadap masing-masing vektor tersebut.
    \begin{answer}
      Determinannya adalah 
      $
        (x_2y_3-x_3y_2)\vec{e}_1
          +(x_3y_1-x_1y_3)\vec{e}_2
          +(x_1y_2-x_2y_1)\vec{e}_3
      $.
      Untuk memeriksa ketegaklurusan, kita periksa bahwa hasil kali titik
      dengan vektor pertama bernilai nol
      \begin{equation*}
        \colvec{x_1 \\ x_2 \\ x_3}
        \dotprod
        \colvec{x_2y_3-x_3y_2 \\ x_3y_1-x_1y_3 \\ x_1y_2-x_2y_1}
        =x_1x_2y_3-x_1x_3y_2+x_2x_3y_1-x_1x_2y_3+x_1x_3y_2-x_2x_3y_1=0
      \end{equation*}
      dan hasil kali titik dengan vektor kedua juga bernilai nol.
      \begin{equation*}
        \colvec{y_1 \\ y_2 \\ y_3}
        \dotprod
        \colvec{x_2y_3-x_3y_2 \\ x_3y_1-x_1y_3 \\ x_1y_2-x_2y_1}
        =x_2y_1y_3-x_3y_1y_2+x_3y_1y_2-x_1y_2y_3+x_1y_2y_3-x_2y_1y_3=0
      \end{equation*}  
    \end{answer}
  \item 
    Buktikan bahwa setiap pernyataan berikut berlaku untuk matriks $\nbyn{2}$.  
    \begin{exparts}
      \partsitem Determinan suatu hasil kali
        merupakan hasil kali determinan-determinannya,
        $\det(ST)=\det(S)\cdot\det(T)$.
      \partsitem Jika \( T \) invertibel, maka
        determinan invers merupakan invers dari determinannya,
        \( \det(T^{-1})=(\,\det(T)\,)^{-1} \).
    \end{exparts}
    Matriks $T$ dan $T^\prime$ disebut 
    \definend{serupa}\index{serupa} jika terdapat matriks 
    nonsingular $P$ sedemikian sehingga $T^\prime=PTP^{-1}$.
    (Kita akan meninjau hubungan ini dalam Bab Lima.)
    Tunjukkan bahwa matriks-matriks \( \nbyn{2} \) yang serupa memiliki
    determinan yang sama.
    \begin{answer}
       \begin{exparts}
        \partsitem Substitusikan dan hitung:
          determinan hasil kalinya adalah
          \begin{align*}
             \det(\begin{mat}
                       a  &b  \\
                       c  &d
                    \end{mat}
                    \begin{mat}
                       w  &x  \\
                       y  &z
                    \end{mat}  )
             &=
             \det(\begin{mat}
                 aw+by  &ax+bz  \\
                 cw+dy  &cx+dz
              \end{mat} )                 \\
             &=
             \begin{array}[t]{@{}l@{}} 
                 acwx+adwz+bcxy+bdyz  \\
                 \> -acwx-bcwz-adxy-bdyz
              \end{array}
          \end{align*}
          sedangkan hasil kali determinannya adalah sebagai berikut.
          \begin{equation*}
             \det(\begin{mat}
                a  &b  \\
                c  &d
             \end{mat})
             \cdot\det(\begin{mat}
                w  &x  \\
                y  &z
             \end{mat})
             =
             (ad-bc)\cdot (wz-xy)
          \end{equation*}
          Mudah untuk memverifikasi bahwa keduanya sama.
        \partsitem Gunakan butir sebelumnya.
       \end{exparts}  
       \noindent Fakta bahwa matriks-matriks serupa memiliki determinan yang sama 
       langsung diperoleh dari kedua hasil di atas:
       $
          \det(PTP^{-1})=\det(P)\cdot\det(T)\cdot\det(P^{-1})
       $.
     \end{answer}
  \recommended \item
    Buktikan bahwa luas daerah berikut pada bidang
    \begin{center}
      \includegraphics{det/mp/ch4.30}
    \end{center}
    sama dengan nilai determinan berikut.
    \begin{equation*}
       \det(
       \begin{mat}
          x_1  &x_2  \\
          y_1  &y_2
       \end{mat})
    \end{equation*}
    Bandingkan dengan determinan berikut.
    \begin{equation*}
       \det(
       \begin{mat}
          x_2  &x_1  \\
          y_2  &y_1
       \end{mat})
    \end{equation*}
    \begin{answer}
      Salah satu caranya adalah menghitung luas-luas berikut
      \begin{center}
        \includegraphics{det/mp/ch4.31}
      \end{center}
      dengan mengambil luas seluruh persegi panjang lalu mengurangkan luas
      $A$, persegi panjang kiri atas; $B$, segitiga tengah atas;
      $D$, segitiga kanan atas; $C$, segitiga kiri bawah; 
      $E$, segitiga tengah bawah; dan $F$, persegi panjang kanan bawah,       
      \( (x_1+x_2)(y_1+y_2)-x_2y_1-(1/2)x_1y_1-(1/2)x_2y_2
              -(1/2)x_2y_2-(1/2)x_1y_1-x_2y_1 \).
      Penyederhanaan memberikan rumus determinan tersebut.

      Determinan ini merupakan negatif dari determinan di atas; rumus tersebut
      membedakan apakah kolom kedua berada berlawanan arah jarum jam terhadap
      kolom pertama.  
     \end{answer}
  \item 
    Buktikan bahwa untuk matriks \( \nbyn{2} \), determinan suatu matriks
    sama dengan determinan transposnya.
    Apakah hal itu juga berlaku untuk matriks \( \nbyn{3} \)?
    \begin{answer}
      Perhitungan untuk matriks \( \nbyn{2} \), dengan menggunakan
      rumus yang dikutip dalam pendahuluan, mudah dilakukan.
      Hal itu juga berlaku untuk matriks \( \nbyn{3} \); 
      perhitungannya rutin.  
    \end{answer}
  \item 
    Apakah fungsi determinan linear\Dash yakni, apakah
    \( \det(x\cdot T+y\cdot S)=x\cdot \det(T)+y\cdot \det(S) \)?
    \begin{answer}
      Tidak.
      Kita menunjukkannya dengan determinan $\nbyn{2}$.
      Ingat bahwa konstanta dikeluarkan satu baris demi satu baris.
      \begin{equation*}
         \det(
         \begin{mat}[r]
            2  &4  \\
            2  &6  \\
         \end{mat})
         =
         2\cdot\det(\begin{mat}[r]
            1  &2  \\
            2  &6  \\
         \end{mat})
         =
         2\cdot 2\cdot \det(\begin{mat}[r]
            1  &2  \\
            1  &3  \\
         \end{mat})
      \end{equation*}
      Hal ini bertentangan dengan linearitas (di sini kita tidak memerlukan \( S \); dengan kata lain, kita dapat 
      mengambil $S$ sebagai matriks nol).  
    \end{answer}
  \item 
    Tunjukkan bahwa jika \( A \) berukuran \( \nbyn{3} \), maka
    \( \det(c\cdot A)=c^3\cdot \det(A) \) untuk setiap skalar \( c \).
    \begin{answer}
       Keluarkan faktor-faktor \( c \) satu baris demi satu baris.  
    \end{answer}
  \item 
    Bilangan real \( \theta \) mana yang membuat
    \begin{equation*}
       \begin{mat}
          \cos\theta  &-\sin\theta  \\
          \sin\theta  &\cos\theta
       \end{mat}
    \end{equation*}
    singular?
    Jelaskan secara geometris.
    \begin{answer}
      Tidak ada bilangan real \( \theta \) yang membuat matriks tersebut singular 
      karena determinan matriks itu,
      \( \cos^2\theta+\sin^2\theta \), tidak pernah bernilai $0$; nilainya sama dengan $1$
      untuk setiap $\theta$.
      Secara geometris, relatif terhadap basis standar,
      matriks ini merepresentasikan
      rotasi bidang sebesar sudut \( \theta \).
      Setiap pemetaan tersebut bersifat satu-ke-satu\Dash antara lain, pemetaan itu invertibel.  
    \end{answer}
  \puzzle \item  
    \cite{Monthly55p257}
    Jika suatu determinan orde tiga memiliki elemen-elemen
    \( 1 \), \( 2 \), \ldots, \( 9 \), berapakah nilai maksimum yang dapat
    dimilikinya?
    \begin{answer}
      \answerasgiven
      Misalkan \( P \) adalah jumlah tiga suku positif determinan tersebut
      dan \( -N \) adalah jumlah tiga suku negatifnya.
      Nilai maksimum \( P \) adalah
      \begin{equation*}
        9\cdot 8\cdot 7 +6\cdot 5\cdot 4 +3\cdot 2\cdot 1=630.
      \end{equation*}
      Nilai minimum \( N \) yang konsisten dengan \( P \) adalah
      \begin{equation*}
        9\cdot 6\cdot 1 +8\cdot 5\cdot 2 +7\cdot 4\cdot 3=218.
      \end{equation*}
      Setiap perubahan pada \( P \) akan menurunkan jumlah tersebut lebih dari
      \( 4 \).
      Karena itu, \( 412 \) merupakan nilai maksimum determinan, dan salah satu bentuk
      determinan yang mencapainya adalah
      \begin{equation*}
         \begin{vmat}[r]
            9  &4  &2  \\
            3  &8  &6  \\
            5  &1  &7
         \end{vmat}.
      \end{equation*}  
    \end{answer}
\end{exercises}




















\subsection{Sifat-Sifat Determinan}
\index{determinan|(}
Kita menginginkan suatu rumus
untuk menentukan apakah matriks $\nbyn{n}$ nonsingular.
Kita tidak akan memulai dengan menyatakan rumus tersebut.
Sebaliknya, untuk setiap~$n$, kita akan mulai dengan meninjau
fungsi yang dihitung oleh rumus semacam itu.
Kita akan mendefinisikan fungsi tersebut melalui daftar sifat.
Kemudian kita akan membuktikan bahwa fungsi dengan sifat-sifat tersebut ada dan tunggal,
serta menjelaskan cara menghitungnya.
(Karena pada akhirnya kita akan membuktikan hal ini, sejak awal 
kita cukup mengatakan `\( \det(T) \)' sebagai pengganti
`jika terdapat fungsi determinan tunggal, maka \( \det(T) \)'.)

\begin{definition}  \label{def:Det}
%<*df:Det>
Suatu \( \nbyn{n} \) \definend{determinan\/}\index{determinan!definisi}%
\index{matriks!determinan} adalah fungsi
\( \map{\det}{\matspace_{\nbyn{n}}}{\Re} \) sedemikian sehingga
\begin{enumerate}
  \item
    $
       \det (\vec{\rho}_1,\dots,k\cdot\vec{\rho}_i 
                                    + \vec{\rho}_j,\dots,\vec{\rho}_n)
       =\det (\vec{\rho}_1,\dots,\vec{\rho}_j,\dots,\vec{\rho}_n)
    $ untuk \( i\ne j \)
  \item 
    $
       \det (\vec{\rho}_1,\ldots,\vec{\rho}_j,
               \dots,\vec{\rho}_i,\dots,\vec{\rho}_n)
       = -\det (\vec{\rho}_1,\dots,\vec{\rho}_i,\dots,\vec{\rho}_j,
                \dots,\vec{\rho}_n)
    $
    untuk \( i\ne j\)
  \item
    $
       \det (\vec{\rho}_1,\dots,k\vec{\rho}_i,\dots,\vec{\rho}_n)
       = k\cdot \det (\vec{\rho}_1,\dots,\vec{\rho}_i,\dots,\vec{\rho}_n)
    $
    untuk setiap skalar~$k$
    % for \( k\ne 0\)
  \item 
     $
       \det(I)=1
     $
     dengan \( I \) suatu matriks identitas
\end{enumerate}
(vektor-vektor $\vec{\rho}\,$ merupakan baris-baris matriks tersebut).
Kita sering menulis \( \deter{T} \) sebagai pengganti \( \det (T) \).
%</df:Det>
\end{definition}

\begin{remark}  \label{rem:SwapRowsRedun}
%<*re:SwapRowsRedun>
Syarat~(2) bersifat redundan karena
 \begin{equation*}
   T\grstep{\rho_i+\rho_j}
    \repeatedgrstep{-\rho_j+\rho_i}
    \repeatedgrstep{\rho_i+\rho_j}
    \repeatedgrstep{-\rho_i}
    \hat{T} 
\end{equation*}
mempertukarkan baris \( i \) dan~\( j \).
Kita mencantumkannya agar 
konsisten dengan penyajian Metode Gauss dalam bab-bab sebelumnya.
%</re:SwapRowsRedun>
\end{remark}

\begin{remark}
Syarat~(3) tidak memiliki pembatasan $k\neq 0$, meskipun
operasi Metode Gauss yang mengalikan suatu baris dengan~$k$
memiliki pembatasan tersebut. 
% \nearbylemma{le:IdenRowsDetZero} below 
Hasil berikut menunjukkan bahwa pembatasan itu tidak diperlukan di sini.
\end{remark}

% % \begin{remark}
% The previous subsection's plan 
% asks for the determinant of an echelon form matrix to be the product 
% down the diagonal.
% The next result shows that in the presence of the other
% three, condition~(4) gives that.    
% % \end{remark}

\begin{lemma}   \label{le:IdenRowsDetZero}
%<*lm:IdenRowsDetZero>
Matriks dengan dua baris identik memiliki determinan nol.
Matriks dengan suatu baris nol memiliki determinan nol.
Suatu matriks nonsingular jika dan hanya jika determinannya tak nol.
Determinan matriks bentuk eselon merupakan hasil kali entri-entri diagonalnya.
%</lm:IdenRowsDetZero>
\end{lemma}

\begin{proof}
%<*pf:IdenRowsDetZero0>
Untuk memverifikasi kalimat pertama, pertukarkan kedua baris yang sama.
Tanda determinan berubah, tetapi matriksnya tetap sama,
sehingga determinannya juga tetap sama.
Jadi, determinannya bernilai nol.
%</pf:IdenRowsDetZero0>

%<*pf:IdenRowsDetZero1>
Untuk kalimat kedua, 
kalikan baris nol dengan dua.
Operasi itu menggandakan determinan, tetapi juga
membiarkan baris tersebut tidak berubah, sehingga 
determinannya tidak berubah. 
Jadi, determinan tersebut harus bernilai nol.
%</pf:IdenRowsDetZero1>

%<*pf:IdenRowsDetZero2>
Untuk kalimat ketiga, lakukan reduksi Gauss--Jordan, 
$T \rightarrow\cdots\rightarrow\hat{T}$. 
Berdasarkan tiga sifat pertama,
determinan $T$ bernilai nol jika dan hanya jika
determinan $\hat{T}$ bernilai nol
(meskipun keduanya dapat berbeda tanda atau besar).
Matriks nonsingular $T$ tereduksi melalui Gauss--Jordan menjadi matriks identitas,
sehingga memiliki determinan tak nol.
Matriks singular $T$ tereduksi menjadi $\hat{T}$ yang memiliki baris nol;
berdasarkan kalimat kedua lema ini, determinannya nol.
%</pf:IdenRowsDetZero2>

%<*pf:IdenRowsDetZero3>
Kalimat keempat memiliki dua kasus.
Jika matriks bentuk eselon tersebut singular, maka
matriks itu memiliki baris nol.
Jadi, diagonalnya memuat nol dan
hasil kali entri-entri diagonalnya bernilai nol.
Berdasarkan kalimat ketiga hasil ini, determinannya nol, sehingga
determinan matriks tersebut sama dengan
hasil kali entri-entri diagonalnya.
%</pf:IdenRowsDetZero3>

%<*pf:IdenRowsDetZero4>
Jika matriks bentuk eselon tersebut nonsingular, maka tidak satu pun entri diagonalnya
bernilai nol.
Ini berarti kita dapat membagi dengan entri-entri tersebut dan 
menggunakan syarat~(3) untuk memperoleh $1$ pada diagonalnya.
\begin{equation*}
  \begin{vmat}
    t_{1,1}  &t_{1,2}  &     &t_{1,n}  \\
    0        &t_{2,2}  &     &t_{2,n}  \\
             &         &\ddots         \\
    0        &         &     &t_{n,n}
  \end{vmat}
  =
  t_{1,1}\cdot t_{2,2}\cdots t_{n,n}\cdot
  \begin{vmat}
    1        &t_{1,2}/t_{1,1}  &     &t_{1,n}/t_{1,1}  \\
    0        &1                &     &t_{2,n}/t_{2,2}  \\
             &                 &\ddots         \\
    0        &                 &     &1
  \end{vmat}
\end{equation*}
Kemudian bagian Jordan dari eliminasi Gauss--Jordan menghasilkan matriks identitas.
\begin{equation*}
  =
  t_{1,1}\cdot t_{2,2}\cdots t_{n,n}\cdot
  \begin{vmat}
    1        &0                &     &0                \\
    0        &1                &     &0                \\
             &                 &\ddots         \\
    0        &                 &     &1
  \end{vmat}
  =
  t_{1,1}\cdot t_{2,2}\cdots t_{n,n}\cdot 1
\end{equation*}
Jadi, dalam kasus ini pun, determinannya
merupakan hasil kali entri-entri diagonal. 
%</pf:IdenRowsDetZero4>
\end{proof}

Hasil tersebut memberi kita cara untuk menghitung nilai fungsi
determinan pada suatu matriks:
lakukan reduksi Gauss sambil mencatat setiap perubahan  
tanda akibat pertukaran baris dan setiap skalar yang kita keluarkan sebagai faktor, 
lalu akhiri dengan mengalikan
entri-entri diagonal hasil berbentuk eselon.
Algoritma ini sama cepatnya dengan Metode Gauss, sehingga
praktis digunakan pada semua matriks yang akan 
kita jumpai.

\begin{example}
Menghitung determinan \( \nbyn{2} \)
dengan Metode Gauss 
\begin{equation*}
   \begin{vmat}[r]
      2  &4  \\
      -1 &3
   \end{vmat}
   =
   \begin{vmat}[r]
      2  &4  \\
      0  &5
   \end{vmat}
   =10
\end{equation*}
tidak banyak menghemat waktu
karena rumus determinan $\nbyn{2}$ mudah digunakan.
Namun, determinan \( \nbyn{3} \) sering kali lebih mudah dihitung
dengan Metode Gauss daripada dengan rumusnya.
\begin{equation*}
   \begin{vmat}[r]
     2  &2  &6  \\
     4  &4  &3  \\
     0  &-3 &5
   \end{vmat}
   =
   \begin{vmat}[r]
     2  &2  &6  \\
     0  &0  &-9 \\
     0  &-3 &5
   \end{vmat}
   =
   -\begin{vmat}[r]
     2  &2  &6  \\
     0  &-3 &5  \\
     0  &0  &-9
   \end{vmat}
   =-54
\end{equation*}
\end{example}

\begin{example}
Determinan yang lebih besar daripada $\nbyn{3}$ dapat dihitung 
dengan cepat menggunakan prosedur Metode Gauss.
\begin{equation*}
   \begin{vmat}[r]
      1  &0  &1  &3  \\
      0  &1  &1  &4  \\
      0  &0  &0  &5  \\
      0  &1  &0  &1
   \end{vmat}
   =
   \begin{vmat}[r]
      1  &0  &1  &3  \\
      0  &1  &1  &4  \\
      0  &0  &0  &5  \\
      0  &0  &-1 &-3
   \end{vmat}
   =
   -\begin{vmat}[r]
      1  &0  &1  &3  \\
      0  &1  &1  &4  \\
      0  &0  &-1 &-3 \\
      0  &0  &0  &5
   \end{vmat}
   =-(-5)=5
\end{equation*}
\end{example}

Contoh tersebut mengangkat suatu hal penting.
Pendahuluan bab ini memberikan rumus untuk determinan $\nbyn{2}$ 
dan $\nbyn{3}$, sehingga kita mengetahui bahwa keduanya ada,
tetapi tidak memberikan rumus fungsi determinan untuk matriks $\nbyn{4}$ atau yang lebih besar.
Sebagai gantinya, \nearbydefinition{def:Det} memberikan sifat-sifat yang 
harus dimiliki suatu fungsi determinan
dan mengarahkan kita untuk menghitung determinan dengan Metode Gauss.

Namun, untuk setiap matriks kita dapat mereduksinya ke bentuk eselon dengan Metode Gauss melalui
berbagai cara.
Sebagai contoh, jika diberikan suatu reduksi, kita dapat mengubahnya dengan menyisipkan langkah pertama
yang mengalikan baris teratas dengan~$2$, lalu 
langkah kedua yang mengalikannya dengan~$1/2$.
Karena itu, kita perlu memastikan bahwa dua reduksi Metode Gauss yang berbeda tidak
menghasilkan dua nilai determinan yang berbeda.

Dengan kata lain, kita harus memverifikasi bahwa \nearbydefinition{def:Det} memberikan 
fungsi yang terdefinisi dengan baik.
Dua subbagian berikutnya melakukan hal tersebut dengan menunjukkan 
bahwa terdapat fungsi terdefinisi baik 
yang memenuhi definisi itu. 

Namun, terlebih dahulu kita tunjukkan bahwa jika fungsi semacam itu ada, 
maka jumlahnya tidak lebih dari satu.
Contoh di atas menggambarkan gagasannya:~kita 
memperoleh~$5$ dengan mengikuti sifat-sifat dalam definisi.
Jadi, meskipun kita belum membuktikan bahwa $\det_{\nbyn{4}}$ ada,
yakni bahwa terdapat fungsi yang memiliki sifat (1)\,--\,(4),  
jika fungsi yang memenuhi sifat-sifat tersebut memang ada, 
kita mengetahui nilai yang diberikannya pada matriks di atas.

\begin{lemma} \label{lm:DetFcnIsUnique}
%<*lm:DetFcnIsUnique>
Untuk setiap $n$, 
jika terdapat fungsi determinan $\nbyn{n}$, maka fungsi tersebut tunggal.
%</lm:DetFcnIsUnique>
\end{lemma}

\begin{proof}
%<*pf:DetFcnIsUnique>
Misalkan terdapat dua fungsi 
$\map{\det_1,\det_2}{\matspace_{\nbyn{n}}}{\Re}$ 
yang memenuhi sifat-sifat dalam 
\nearbydefinition{def:Det} beserta konsekuensinya, \nearbylemma{le:IdenRowsDetZero}.
Untuk suatu matriks persegi~$M$, tetapkan salah satu cara melakukan Metode Gauss
untuk membawa matriks itu ke bentuk eselon (adanya
berbagai cara tidak menjadi masalah; cukup tetapkan salah satunya).
Dengan menggunakan reduksi tetap ini seperti pada
contoh-contoh di atas\Dash 
mencatat faktor penskalaan baris dan pergantian tanda pada pertukaran baris, 
lalu mengalikan entri-entri diagonal hasil berbentuk eselon\Dash
kita dapat menghitung nilai yang harus dikembalikan kedua fungsi tersebut pada~$M$,
dan keduanya harus mengembalikan nilai yang sama.
Karena keduanya memberikan keluaran yang sama untuk setiap masukan, keduanya merupakan fungsi yang sama.
%</pf:DetFcnIsUnique>
\end{proof}

Frasa `jika terdapat fungsi determinan $\nbyn{n}$' 
menekankan bahwa,
meskipun kita dapat
menggunakan Metode Gauss untuk menghitung satu-satunya nilai yang mungkin 
dikembalikan oleh fungsi determinan, 
kita belum menunjukkan bahwa fungsi semacam itu ada untuk semua $n$.
Sisa bagian ini akan menunjukkannya.



\begin{exercises}
  \item[{\em Untuk soal-soal berikut, anggap bahwa fungsi determinan $\nbyn{n}$
      ada untuk semua $n$.}]
  \recommended \item Tentukan setiap determinan dengan melakukan satu operasi baris.
    \begin{exparts*}
      \partsitem \( \begin{vmat}[r]
                 1  &-2  &1  &2 \\
                 2  &-4  &1  &0 \\
                 0  &0  &-1  &0 \\
                 0  &0  &0  &5
               \end{vmat} \)
      \partsitem \( \begin{vmat}[r]
                 1  &1  &-2  \\
                 0  &0  &4  \\
                 0  &3  &-6
               \end{vmat} \)
    \end{exparts*}
    \begin{answer}
      \begin{exparts}
       \partsitem Lakukan $2\rho_1+\rho_2$ untuk memperoleh bentuk eselon, lalu
          kalikan entri-entri diagonalnya.
          Determinannya adalah~$0$.
          % sage: M = matrix(QQ, [[1,-2,1,2], [2,-4,1,0], [0,0,-1,0], [0,0,0,5]])
          % sage: M.determinant()
          % 0
        \partsitem Mempertukarkan baris kedua dan ketiga membawa sistem ke
          bentuk eselon (dan mengubah tanda determinan).
          Mengalikan entri-entri diagonal memberikan~$12$, sehingga determinan 
          matriks yang diberikan adalah~$-12$.
          % sage: M = matrix(QQ, [[1,1,-2], [0,0,4], [0,3,-6]])
          % sage: M.determinant()
          % -12
      \end{exparts}
    \end{answer}
  \recommended \item 
    Gunakan Metode Gauss untuk menentukan setiap determinan.
    \begin{exparts*}
      \partsitem \( \begin{vmat}[r]
                 3  &1  &2  \\
                 3  &1  &0  \\
                 0  &1  &4
               \end{vmat} \)
      \partsitem \( \begin{vmat}[r]
                 1  &0  &0  &1 \\
                 2  &1  &1  &0 \\
                -1  &0  &1  &0 \\
                 1  &1  &1  &0
               \end{vmat} \)
    \end{exparts*}
    \begin{answer}
      \begin{exparts}
      \partsitem \( \begin{vmat}[r]
                 3  &1  &2  \\
                 3  &1  &0  \\
                 0  &1  &4
               \end{vmat}
               =
               \begin{vmat}[r]
                 3  &1  &2  \\
                 0  &0  &-2 \\
                 0  &1  &4
               \end{vmat}
               =
               -\begin{vmat}[r]
                 3  &1  &2  \\
                 0  &1  &4  \\
                 0  &0  &-2
               \end{vmat}
               =6  \)
      \partsitem \( \begin{vmat}[r]
                 1  &0  &0  &1 \\
                 2  &1  &1  &0 \\
                -1  &0  &1  &0 \\
                 1  &1  &1  &0
               \end{vmat}
               =
               \begin{vmat}[r]
                 1  &0  &0  &1 \\
                 0  &1  &1  &-2\\
                 0  &0  &1  &1 \\
                 0  &1  &1  &-1
               \end{vmat}
               =
               \begin{vmat}[r]
                 1  &0  &0  &1 \\
                 0  &1  &1  &-2\\
                 0  &0  &1  &1 \\
                 0  &0  &0  &1
               \end{vmat}
               =1    \)
      \end{exparts} 
    \end{answer}
  \item Gunakan Metode Gauss untuk menentukan masing-masing determinan berikut.
    \begin{exparts*}
      \partsitem \( \begin{vmat}[r]
                 2  &-1  \\
                 -1 &-1
               \end{vmat}  \)
      \partsitem \( \begin{vmat}[r]
                 1  &1  &0  \\
                 3  &0  &2  \\
                 5  &2  &2
               \end{vmat} \)
    \end{exparts*}
    \begin{answer}
     \begin{exparts}
      \partsitem \(
        \begin{vmat}[r]
                 2  &-1  \\
                 -1 &-1
               \end{vmat}
             =\begin{vmat}[r]
                 2  &-1  \\
                 0  &-3/2
               \end{vmat}=-3  \);
      \partsitem \( \begin{vmat}[r]
                 1  &1  &0  \\
                 3  &0  &2  \\
                 5  &2  &2
               \end{vmat}
              =\begin{vmat}[r]
                 1  &1  &0  \\
                 0  &-3 &2  \\
                 0  &-3 &2
               \end{vmat}
              =\begin{vmat}[r]
                 1  &1  &0  \\
                 0  &-3 &2  \\
                 0  &0  &0
               \end{vmat}
              =0 \)
     \end{exparts}  
    \end{answer}
  \item 
    Untuk nilai \( k \) yang mana sistem berikut 
    memiliki solusi tunggal?
    \begin{equation*}
       \begin{linsys}{4}
          x  &  &  &+ &z  &-  &w  &=  &2  \\
             &  &y &- &2z &   &   &=  &3  \\
          x  &  &  &+ &kz &   &   &=  &4  \\
             &  &  &  &z  &-  &w  &=  &2
        \end{linsys}
    \end{equation*}
    \begin{answer}
      Kapan determinannya bernilai tak nol?
      \begin{equation*}
         \begin{vmat}
           1  &0  &1  &-1  \\
           0  &1  &-2 &0   \\
           1  &0  &k  &0   \\
           0  &0  &1  &-1
         \end{vmat}
         =
         \begin{vmat}
           1  &0  &1  &-1  \\
           0  &1  &-2 &0   \\
           0  &0  &k-1&1   \\
           0  &0  &1  &-1
         \end{vmat}           
      \end{equation*}
      Jelas bahwa 
      $k=1$ menghasilkan kenonsingularan, sehingga determinannya tak nol. 
      Jika $k\neq 1$, maka 
      kita memperoleh bentuk eselon dengan kombinasi $(-1/(k-1))\rho_3+\rho_4$.
      \begin{equation*}
         =
         \begin{vmat}
           1  &0  &1  &-1  \\
           0  &1  &-2 &0   \\
           0  &0  &k-1&1   \\
           0  &0  &0  &-1-1/(k-1)
         \end{vmat}
      \end{equation*}
      Mengalikan entri-entri diagonal memberikan $(k-1)(-1-1/(k-1))=-(k-1)-1=-k$.
      Jadi, matriks tersebut memiliki determinan tak nol, dan dengan demikian sistemnya
      memiliki solusi tunggal, jika dan hanya jika \( k\neq 0 \).  
    \end{answer}
  \recommended \item 
    Nyatakan masing-masing bentuk berikut dalam \( \deter{H} \).
    \begin{exparts}
      \partsitem \( \begin{vmat}
            h_{3,1}  &h_{3,2} &h_{3,3} \\
            h_{2,1}  &h_{2,2} &h_{2,3} \\
            h_{1,1}  &h_{1,2} &h_{1,3}
          \end{vmat} \)
      \partsitem \( \begin{vmat}
           -h_{1,1}   &-h_{1,2}  &-h_{1,3} \\
           -2h_{2,1}  &-2h_{2,2} &-2h_{2,3} \\
           -3h_{3,1}  &-3h_{3,2} &-3h_{3,3}
          \end{vmat} \)
      \partsitem \( \begin{vmat}
            h_{1,1}+h_{3,1}  &h_{1,2}+h_{3,2} &h_{1,3}+h_{3,3} \\
            h_{2,1}          &h_{2,2}         &h_{2,3} \\
            5h_{3,1}         &5h_{3,2}        &5h_{3,3}
          \end{vmat} \)
    \end{exparts}
    \begin{answer}
      \begin{exparts}
        \partsitem Syarat~(2) dalam definisi determinan
          berlaku melalui pertukaran $\rho_1\leftrightarrow\rho_3$.
          \begin{equation*}
            \begin{vmat}
              h_{3,1}  &h_{3,2} &h_{3,3} \\
              h_{2,1}  &h_{2,2} &h_{2,3} \\
              h_{1,1}  &h_{1,2} &h_{1,3}
            \end{vmat}
            =
            -\begin{vmat}
              h_{1,1}  &h_{1,2} &h_{1,3} \\
              h_{2,1}  &h_{2,2} &h_{2,3} \\
              h_{3,1}  &h_{3,2} &h_{3,3} 
            \end{vmat}
          \end{equation*}
        \partsitem Syarat~(3) berlaku.
          \begin{multline*}
            \begin{vmat}
             -h_{1,1}   &-h_{1,2}  &-h_{1,3} \\
             -2h_{2,1}  &-2h_{2,2} &-2h_{2,3} \\
             -3h_{3,1}  &-3h_{3,2} &-3h_{3,3}
            \end{vmat}
            =
            (-1)\cdot(-2)\cdot(-3)\cdot
            \begin{vmat}
             h_{1,1}   &h_{1,2}  &h_{1,3} \\
             h_{2,1}   &h_{2,2}  &h_{2,3} \\
             h_{3,1}   &h_{3,2}  &h_{3,3}
            \end{vmat}                                  \\
            =
            (-6)\cdot
            \begin{vmat}
             h_{1,1}   &h_{1,2}  &h_{1,3} \\
             h_{2,1}   &h_{2,2}  &h_{2,3} \\
             h_{3,1}   &h_{3,2}  &h_{3,3}
            \end{vmat}
          \end{multline*}
        \partsitem 
          \begin{multline*}
          \begin{vmat}
            h_{1,1}+h_{3,1}  &h_{1,2}+h_{3,2} &h_{1,3}+h_{3,3} \\
            h_{2,1}          &h_{2,2}         &h_{2,3} \\
            5h_{3,1}         &5h_{3,2}        &5h_{3,3}
          \end{vmat}                                                 \\
          \begin{aligned}
          &=
          5\cdot     
          \begin{vmat}
            h_{1,1}+h_{3,1}  &h_{1,2}+h_{3,2} &h_{1,3}+h_{3,3} \\
            h_{2,1}          &h_{2,2}         &h_{2,3} \\
            h_{3,1}          &h_{3,2}         &h_{3,3}
          \end{vmat}                                     \\
          &=5\cdot
          \begin{vmat}
            h_{1,1}          &h_{1,2}         &h_{1,3} \\
            h_{2,1}          &h_{2,2}         &h_{2,3} \\
            h_{3,1}          &h_{3,2}         &h_{3,3}
          \end{vmat}
          \end{aligned}
          \end{multline*}
      \end{exparts}  
     \end{answer}
  \recommended \item 
    Tentukan determinan suatu matriks diagonal.
    \begin{answer}
       Matriks diagonal berada dalam bentuk eselon, sehingga
       determinannya merupakan hasil kali entri-entri diagonalnya.
    \end{answer}
  \item 
    Jelaskan himpunan solusi suatu sistem linear homogen jika
    determinan matriks koefisiennya tak nol.
    \begin{answer}
      Himpunan tersebut adalah subruang trivial.  
    \end{answer}
  \recommended \item 
    Tunjukkan bahwa determinan berikut bernilai nol.
    \begin{equation*}
      \begin{vmat}
        y+z  &x+z  &x+y  \\
        x    &y    &z    \\
        1    &1    &1
      \end{vmat}
    \end{equation*}
    \begin{answer} 
      Menambahkan baris kedua ke baris pertama memberikan matriks yang baris pertamanya
      sama dengan $x+y+z$ kali baris ketiganya.  
    \end{answer}
  \item 
    \begin{exparts}
      \partsitem Tentukan matriks $\nbyn{1}$, $\nbyn{2}$, dan $\nbyn{3}$
        yang entri $i,j$-nya diberikan oleh $(-1)^{i+j}$.  
      \partsitem Tentukan determinan matriks persegi yang entri \( i,j \)-nya
        adalah \( (-1)^{i+j} \).
    \end{exparts}
    \begin{answer}
      \begin{exparts}
        \partsitem 
          $\begin{mat}[r]
            1
          \end{mat}$,
          $\begin{mat}[r]
            1  &-1  \\
            -1  &1
          \end{mat}$,
          $\begin{mat}[r]
            1   &-1  &1   \\
            -1  &1   &-1  \\
            1   &-1  &1
          \end{mat}$
        \partsitem Determinan pada kasus $\nbyn{1}$ adalah $1$.
          Pada setiap kasus lainnya, baris kedua merupakan negatif dari baris pertama,
          sehingga matriks tersebut singular dan determinannya nol. 
      \end{exparts} 
    \end{answer}
  \item 
    \begin{exparts}
      \partsitem Tentukan matriks $\nbyn{1}$, $\nbyn{2}$, dan $\nbyn{3}$
        yang entri $i,j$-nya diberikan oleh $i+j$.  
      \partsitem Tentukan determinan matriks persegi yang 
        entri $i,j$-nya adalah $i+j$.
    \end{exparts}
    \begin{answer}
      \begin{exparts}
        \partsitem 
          $\begin{mat}[r]
            2
          \end{mat}$,
          $\begin{mat}[r]
            2  &3  \\
            3  &4
          \end{mat}$,
          $\begin{mat}[r]
            2  &3  &4  \\
            3  &4  &5  \\
            4  &5  &6
          \end{mat}$
        \partsitem Kasus $\nbyn{1}$ dan $\nbyn{2}$ memberikan hasil berikut.
          \begin{equation*}
            \begin{vmat}[r]
              2
            \end{vmat}
            =2
            \qquad
            \begin{vmat}[r]
              2  &3  \\
              3  &4
            \end{vmat}=-1
          \end{equation*}
          Matriks $\nbyn{n}$ dengan $n\geq 3$ bersifat singular; sebagai contoh,
          \begin{equation*}
            \begin{vmat}[r]
              2  &3  &4  \\
              3  &4  &5  \\
              4  &5  &6
            \end{vmat}=0
          \end{equation*}
          karena dua kali baris kedua dikurangi baris pertama 
          sama dengan baris ketiga.
          Pemeriksaan ini bersifat rutin.
      \end{exparts} 
    \end{answer}
  \recommended \item 
    Tunjukkan bahwa fungsi determinan tidak linear dengan  
    memberikan kasus ketika \( \deter{A+B}\neq\deter{A}+\deter{B} \).
    \begin{answer}
      Contoh berikut 
      \begin{equation*}
         A=B=
         \begin{mat}[r]
           1  &2  \\
           3  &4
         \end{mat}
      \end{equation*}
      mudah diperiksa.
      \begin{equation*}
         \deter{A+B}
         =
         \begin{vmat}[r]
           2  &4  \\
           6  &8
         \end{vmat}
         =-8
         \qquad
         \deter{A}+\deter{B}
         =
         -2-2
         =-4
      \end{equation*}
      Kebetulan, contoh ini juga memberikan kasus ketika perkalian skalar 
      tidak dipertahankan, 
      $\deter{2\cdot A}\neq 2\cdot\deter{A}$.  
     \end{answer}
  \item 
    Syarat kedua dalam definisi, yaitu bahwa pertukaran baris mengubah 
    tanda determinan, agak merepotkan.
    Artinya, kita harus mencatat banyaknya pertukaran untuk menghitung bagaimana
    tandanya berganti-ganti.
    Dapatkah kita menyingkirkannya?
    Dapatkah kita menggantinya dengan syarat bahwa pertukaran baris membiarkan determinan
    tidak berubah?
    (Jika demikian, kita akan memerlukan rumus baru untuk $\nbyn{1}$,
     $\nbyn{2}$, dan $\nbyn{3}$, tetapi itu hanya persoalan kecil.)
     \begin{answer}
       Tidak, kita tidak dapat menggantinya.
       \nearbyremark{rem:SwapRowsRedun}
       menunjukkan bahwa keempat syarat setelah penggantian tersebut akan  
       saling bertentangan\Dash tidak ada fungsi yang memenuhi semuanya.
     \end{answer}
  \item 
    Buktikan bahwa determinan setiap matriks segitiga, baik atas maupun bawah,
    merupakan hasil kali entri-entri diagonalnya.
    \begin{answer}
      Matriks segitiga atas berada dalam bentuk eselon.

      Matriks segitiga bawah bersifat singular atau nonsingular.
      Jika singular, matriks itu memiliki nol pada diagonalnya, sehingga 
      determinannya (yakni nol) memang merupakan hasil kali entri-entri diagonal. 
      Jika nonsingular, matriks itu tidak memiliki nol pada diagonalnya, dan
      kita dapat mereduksinya dengan Metode Gauss ke
      bentuk eselon tanpa mengubah diagonalnya.  
    \end{answer}
  \item
    Rujuklah definisi matriks elementer dalam subbagian Mekanika
    Perkalian Matriks.
    \begin{exparts}
      \partsitem Berapakah determinan setiap jenis matriks elementer?
      \partsitem Buktikan bahwa jika \( E \) adalah suatu matriks elementer, maka
        \( \deter{ES}=\deter{E}\deter{S} \) untuk setiap \( S \)
        yang berukuran sesuai.
      \partsitem \textit{(Pertanyaan ini tidak melibatkan determinan.)}
        Buktikan bahwa jika \( T \) singular, maka hasil kali \( TS \)
        juga singular.
      \partsitem Tunjukkan bahwa \( \deter{TS}=\deter{T}\deter{S} \).
      \partsitem Tunjukkan bahwa jika \( T \) nonsingular, maka
          \( \deter{T^{-1}}=\deter{T}^{-1} \).
    \end{exparts}
    \begin{answer}
      \begin{exparts}
        \partsitem Sifat-sifat dalam definisi determinan 
          menunjukkan bahwa 
          \( \deter{M_i(k)}=k \),
          \( \deter{P_{i,j}}=-1 \),
          dan
          \( \deter{C_{i,j}(k)}=1 \).
        \partsitem Ketiga kasus mudah diperiksa dengan mengingat aksi
          perkalian dari kiri oleh setiap jenis matriks.
        \partsitem Jika \( TS \) invertibel, \( (TS)M=I \), maka
          sifat asosiatif perkalian matriks, 
          \( T(SM)=I \), menunjukkan bahwa \( T \) invertibel.
          Jadi, jika \( T \) tidak invertibel, \( TS \) juga tidak invertibel.
        \partsitem Jika \( T \) singular, terapkan jawaban sebelumnya:
          \( \deter{TS}=0 \) dan
          \( \deter{T}\cdot\deter{S}=0\cdot\deter{S}=0 \).
          Jika \( T \) tidak singular, kita dapat menuliskannya sebagai hasil kali
          matriks-matriks elementer,
          $
            \deter{TS}=\deter{E_r\cdots E_1S}=
            \deter{E_r}\cdots\deter{E_1}\cdot\deter{S}=
            \deter{E_r\cdots E_1}\deter{S}=\deter{T}\deter{S}
          $.
        \partsitem \( 1=\deter{I}=\deter{T\cdot T^{-1}}
                     =\deter{T}\deter{T^{-1}} \)
      \end{exparts}  
     \end{answer} 
  \item 
    Buktikan bahwa determinan suatu hasil kali merupakan hasil kali
    determinannya, \( \deter{TS}=\deter{T}\,\deter{S} \), dengan cara berikut.
    Tetapkan matriks \( \nbyn{n} \) \( S \), lalu tinjau fungsi
    \( \map{d}{\matspace_{\nbyn{n}}}{\Re} \) yang diberikan oleh
    \( T\mapsto \deter{TS}/\deter{S} \).
    \begin{exparts}
      \partsitem Periksa bahwa \( d \) memenuhi syarat~(1) dalam definisi
        fungsi determinan.
      \partsitem Periksa syarat~(2).
      \partsitem Periksa syarat~(3).
      \partsitem Periksa syarat~(4).
      \partsitem Simpulkan bahwa determinan suatu hasil kali merupakan hasil kali
        determinannya.
    \end{exparts}
    \begin{answer}
      \begin{exparts}
        \partsitem Kita harus menunjukkan bahwa jika
          \begin{equation*}
            T\grstep{k\rho_i+\rho_j}\hat{T}
          \end{equation*}
          maka $d(T)=\deter{TS}/\deter{S}
          =\deter{\hat{T}S}/\deter{S}=d(\hat{T})$.
          Bukti selesai jika kita menunjukkan bahwa terlebih dahulu mengombinasikan baris,
          lalu mengalikan untuk memperoleh \( \hat{T}S \), memberikan hasil yang sama dengan
          terlebih dahulu mengalikan untuk memperoleh \( TS \), lalu melakukan kombinasi
          (karena determinan \( \deter{TS} \) tidak berubah oleh
          kombinasi tersebut, sehingga kita memperoleh \( \deter{\hat{T}S}=\deter{TS} \), dan
          dengan demikian \( d(\hat{T})=d(T) \)).
          Argumennya berjalan sebagai berikut:~setelah menambahkan 
          \( k \) kali baris~\( i \) dari \( TS \) ke
          baris~$j$ dari \( TS \), entri \( j,p \)-nya adalah
          \( (kt_{i,1}+t_{j,1})s_{1,p}+\dots+(kt_{i,r}+t_{j,r})s_{r,p} \),
          yang merupakan entri \( j,p \) dari \( \hat{T}S \).
        \partsitem Kita hanya perlu menunjukkan bahwa melakukan pertukaran
          $
            T\smash[b]{\grstep{\rho_i\swap\rho_j}}\hat{T}
          $
          lalu mengalikan untuk memperoleh \( \hat{T}S \) memberikan hasil yang sama dengan
          mengalikan \( T \) dengan \( S \), lalu melakukan pertukaran (karena
          determinan \( \deter{TS} \) berganti tanda pada
          pertukaran baris, kita kemudian memperoleh \( \deter{\hat{T}S}=-\deter{TS} \),
          sehingga \( d(\hat{T})=-d(T) \)).
          Argumen tersebut berjalan persis seperti argumen sebelumnya.
        \partsitem Kini tidak mengejutkan bahwa kita hanya perlu menunjukkan bahwa 
          mengalikan suatu baris dengan skalar
          $
            T\smash[b]{\grstep{k\rho_i}}\hat{T}
          $
          lalu menghitung \( \hat{T}S \) memberikan hasil yang sama dengan
          terlebih dahulu menghitung \( TS \), lalu mengalikan barisnya dengan \( k \)
          (karena determinan \( \deter{TS} \) ikut terskala dengan \( k \)
          akibat perkalian tersebut, kita memperoleh \( \deter{\hat{T}S}=k\deter{TS} \),
          sehingga \( d(\hat{T})=k\,d(T) \)).
          Argumennya berjalan persis seperti di atas.
        \partsitem Jelas.
        \partsitem Karena kita telah menunjukkan bahwa \( d(T) \) merupakan determinan
          dan bahwa fungsi determinan (jika ada)
          bersifat tunggal, kita memperoleh
          \( \deter{T}=d(T)=\deter{TS}/\deter{S} \).
      \end{exparts}  
    \end{answer}
  \item 
    Suatu \definend{submatriks}\index{submatriks}\index{matriks!submatriks} 
dari matriks $A$ adalah matriks yang diperoleh dengan menghapus 
    beberapa baris dan kolom $A$.
    Jadi, matriks pertama berikut merupakan submatriks dari matriks kedua.
    \begin{equation*}
      \begin{mat}[r]
        3  &1  \\
        2  &5
      \end{mat}
      \qquad
      \begin{mat}[r]
        3  &4  &1  \\
        0  &9  &-2 \\
        2  &-1 &5
      \end{mat}
    \end{equation*}
    Buktikan bahwa untuk setiap matriks persegi,
    rank matriks tersebut adalah $r$ jika dan hanya jika
    \( r \) merupakan bilangan bulat terbesar
    sedemikian sehingga terdapat submatriks \( \nbyn{r} \) dengan
    determinan tak nol.
    \begin{answer}
      Pertama-tama kita akan berargumen bahwa matriks berank \( r \) memiliki
      submatriks \( \nbyn{r} \) dengan determinan tak nol.
      Matriks berank \( r \) memiliki himpunan \( r \) baris yang bebas linear.
      Matriks yang dibentuk dari baris-baris tersebut akan memiliki rank baris \( r \), sehingga
      rank kolomnya juga \( r \).
      Kesimpulannya: dari \( r \) baris tersebut kita dapat mengambil
      himpunan \( r \) kolom yang bebas linear, sehingga matriks semula memiliki
      submatriks \( \nbyn{r} \) berank \( r \).

      Kita menutup dengan menunjukkan bahwa jika \( r \) merupakan bilangan bulat terbesar semacam itu,
      maka rank matriks tersebut adalah \( r \).
      Berdasarkan kemaksimalan \( r \), kita hanya perlu menunjukkan
      bahwa jika suatu matriks memiliki submatriks \( \nbyn{k} \)
      berdeterminan tak nol, maka rank matriks tersebut sekurang-kurangnya \( k \).
      Tinjaulah submatriks \( \nbyn{k} \) semacam itu.
      Baris-barisnya merupakan bagian dari baris-baris matriks semula; jelas bahwa
      himpunan baris utuh yang bersesuaian bebas linear.
      Jadi, rank baris matriks semula sekurang-kurangnya \( k \), dan rank
      baris suatu matriks sama dengan ranknya.  
   \end{answer}
  \item
    Buktikan bahwa matriks dengan entri-entri rasional memiliki determinan rasional.
    \begin{answer}
      Matriks yang hanya memiliki entri rasional tereduksi dengan
      Metode Gauss menjadi matriks bentuk eselon hanya menggunakan aritmetika rasional.
      Jadi, entri-entri diagonalnya harus rasional, sehingga hasil kali
      entri-entri diagonal itu rasional.  
    \end{answer}
  \puzzle \item 
     \cite{Monthly53p115}
     Temukan unsur kesamaan dalam (a)~menyederhanakan pecahan, (b)~membubuhkan bedak
     pada hidung, (c)~membangun anak tangga baru pada gereja, (d)~mempertahankan profesor
     emeritus di kampus, dan (e)~menempatkan \( B \), \( C \), \( D \) dalam
     determinan
     \begin{equation*}
       \begin{vmat}
         1   &a   &a^2  &a^3  \\
         a^3 &1   &a    &a^2  \\
         B   &a^3 &1    &a    \\
         C   &D   &a^3  &1
       \end{vmat}.
     \end{equation*}
     \begin{answer}
       \answerasgiven
       Nilai determinan \( (1-a^4)^3 \) tidak bergantung pada
       nilai \( B \), \( C \), maupun \( D \).
       Karena itu, operasi~(e) tidak mengubah nilai determinan,
       melainkan hanya mengubah tampilannya.
       Jadi, unsur kesamaan dalam (a), (b), (c), (d), dan (e)
       hanyalah bahwa tampilan entitas utamanya berubah.
       Unsur yang sama muncul dalam (f)~mengganti label nama sekuntum mawar,
       (g)~menulis bilangan bulat desimal dalam basis \( 12 \), (h)~menyepuh
       bunga lili, (i)~memoles citra seorang politikus, dan (j)~menganugerahkan gelar
       kehormatan.  
    \end{answer}
\end{exercises}



















\subsection{Ekspansi Permutasi}
Subbagian sebelumnya mendefinisikan suatu fungsi sebagai determinan jika fungsi tersebut
memenuhi empat syarat dan
menunjukkan bahwa terdapat paling banyak satu fungsi determinan $\nbyn{n}$ untuk
setiap $n$.
Yang tersisa adalah menunjukkan bahwa untuk setiap $n$, fungsi semacam itu ada.

Namun, kita dapat menghitung determinan dengan mudah:
kita menggunakan Metode Gauss, mencatat perubahan tanda akibat pertukaran baris, 
lalu mengalikan entri-entri diagonal pada akhir proses.
Bagaimana mungkin fungsi determinan tidak ada?

Kesulitannya ialah menunjukkan bahwa 
perhitungan tersebut memberikan hasil yang terdefinisi dengan baik\Dash yakni, tunggal\Dash.
%<*DifferentGaussMethodReductions>
Tinjaulah dua reduksi Metode Gauss berikut
terhadap matriks yang sama; reduksi pertama tidak menggunakan pertukaran baris,
\begin{equation*}
  \begin{mat}[r]
    1  &2  \\
    3  &4
  \end{mat}
  \grstep{-3\rho_1+\rho_2}
  \begin{mat}[r]
    1  &2  \\
    0  &-2
  \end{mat}
\end{equation*}
sedangkan reduksi kedua menggunakannya sekali.
\begin{equation*}
  \begin{mat}[r]
    1  &2  \\
    3  &4
  \end{mat}
  \grstep{\rho_1\leftrightarrow\rho_2}
  \begin{mat}[r]
    3  &4  \\
    1  &2
  \end{mat}
  \grstep{-(1/3)\rho_1+\rho_2}
  \begin{mat}[r]
    3  &4  \\
    0  &2/3
  \end{mat}
\end{equation*}
Keduanya menghasilkan determinan $-2$ 
karena pada reduksi kedua
kita memperhitungkan bahwa pertukaran baris mengubah tanda hasil
yang diperoleh dengan mengalikan entri-entri diagonal.
%</DifferentGaussMethodReductions>
Fakta bahwa kita dapat melanjutkan melalui dua cara membuka kemungkinan bahwa
keduanya memberikan jawaban yang berbeda.
% To illustrate how a computation that is like the ones that we are doing could 
% fail to be well-defined, 
% suppose that \nearbydefinition{def:Det} did not include condition~(2). 
% That is, suppose that we instead tried to define determinants so that
% the value would not change on a row swap.
% Then first reduction above would
% yield $-2$ while the second would yield $+2$.
% We could still do computations but they wouldn't give consistent outcomes\Dash
% there is no 
% function that satisfies conditions (1), (3), (4), and also this
% altered second condition. 
% Of course, observing that \nearbydefinition{def:Det} does the right thing
% with these two reductions of the above matrix is not enough.
Dengan kata lain, cara yang telah kita berikan untuk menghitung nilai determinan tidak
serta-merta meniadakan kemungkinan bahwa, misalnya,
terdapat dua reduksi suatu matriks $\nbyn{7}$ yang menghasilkan 
nilai determinan berbeda.
Dalam kasus tersebut kita tidak akan memiliki fungsi, karena menurut definisinya,
setiap masukan suatu fungsi harus dipasangkan dengan tepat satu keluaran. 
Sisa bagian ini 
menunjukkan bahwa definisi~\nearbydefinition{def:Det} 
tidak pernah menimbulkan konflik.

% \begin{remark}
% The natural way approach to this would be 
% to define the determinant function with:~``The value of the
% function is the result of doing Gauss's Method,
% keeping track of row swaps, and finishing by multiplying down the diagonal''. 
% (Since Gauss's Method allows for some variation, as we saw above, 
% we would have to fix an explicit algorithm.)
% Then we would be done if we verified that 
% this way algorithm satisfies the four properties.
% However, how to verify this is not evident.
% So the development below will not proceed in this way.
% \end{remark}

Untuk melakukannya, kita akan mendefinisikan cara alternatif untuk menentukan
nilai suatu determinan.
(Dalam praktiknya alternatif ini kurang berguna karena 
lambat.
Namun, alternatif ini sangat berguna untuk teori.)
Gagasan kuncinya ialah bahwa   
syarat~(3) dari \nearbydefinition{def:Det}
menunjukkan bahwa fungsi determinan tidak linear.

\begin{example}
Berdasarkan syarat~(3), skalar dikeluarkan secara terpisah dari setiap baris,
\begin{equation*}
   \begin{vmat}[r]
       4  &2  \\
      -2  &6
     \end{vmat}
  =2\cdot\begin{vmat}[r]
       2  &1  \\
      -2  &6
     \end{vmat}
  =4\cdot\begin{vmat}[r]
       2  &1  \\
      -1  &3
     \end{vmat}
\end{equation*}
bukan dari seluruh matriks sekaligus.
Jadi, jika
\begin{equation*}
   A=\begin{mat}[r]
       2  &1  \\
      -1  &3
     \end{mat}
\end{equation*} 
maka \( \det(2A) \neq 2\cdot\det(A) \) 
(melainkan, \( \det(2A) =  4\cdot\det(A) \)).
\end{example}

Karena skalar dikeluarkan satu baris demi satu baris, kita mungkin menduga bahwa
determinan bersifat linear pada setiap baris secara terpisah.

\begin{definition} \label{def:multilinear}
%<*df:Multilinear>
Misalkan \( V \) suatu ruang vektor.
Suatu pemetaan \( \map{f}{V^n}{\Re} \) disebut
\definend{multilinear}\index{fungsi!multilinear}\index{multilinear}
jika  
\begin{enumerate}
  \item 
    $
      f(\vec{\rho}_1,\dots,\vec{v}+\vec{w},
      \ldots,\vec{\rho}_n)
      =f(\vec{\rho}_1,\dots,\vec{v},\dots,\vec{\rho}_n)
      +f(\vec{\rho}_1,\dots,\vec{w},\dots,\vec{\rho}_n)
    $
  \item 
    $
      f(\vec{\rho}_1,\dots,k\vec{v},\dots,\vec{\rho}_n)
      =k\cdot f(\vec{\rho}_1,\dots,\vec{v},\dots,\vec{\rho}_n)
    $
\end{enumerate}
untuk \( \vec{v}, \vec{w}\in V \) dan \( k\in\Re \).
%</df:Multilinear>
\end{definition}

\begin{lemma}  \label{lem:DetsMultilinear}
%<*lm:DetsMultilinear>
Determinan bersifat multilinear.
%</lm:DetsMultilinear>
\end{lemma}

\begin{proof}
%<*pf:DetsMultilinear0>
Sifat~(2) di sini tepat merupakan syarat~(3) dalam \nearbydefinition{def:Det},
sehingga kita hanya perlu memverifikasi sifat~(1).
%</pf:DetsMultilinear0>
% \begin{equation*}
%   \det(\vec{\rho}_1,\dots,\vec{v}+\vec{w},
%       \dots,\vec{\rho}_n)
%   =\det(\vec{\rho}_1,\dots,\vec{v},\dots,\vec{\rho}_n)
%   +\det(\vec{\rho}_1,\dots,\vec{w},\dots,\vec{\rho}_n)
% \end{equation*}

%<*pf:DetsMultilinear1>
Terdapat dua kasus.
Jika himpunan baris lainnya
\(
  \set{\vec{\rho}_1,\dots,\vec{\rho}_{i-1},\vec{\rho}_{i+1},
       \dots,\vec{\rho}_n} 
\)
bergantung linear, maka ketiga matriks tersebut singular, sehingga ketiga
determinannya nol dan kesamaannya langsung berlaku.
%</pf:DetsMultilinear1>

%<*pf:DetsMultilinear2>
Karena itu, anggap bahwa himpunan baris lainnya bebas linear.
Kita dapat membentuk basis dengan 
menambahkan satu vektor lagi,
$\sequence{\vec{\rho}_1,\dots,\vec{\rho}_{i-1},\vec{\beta},
               \vec{\rho}_{i+1},\dots,\vec{\rho}_n}$.
Nyatakan $\vec{v}$ dan $\vec{w}$ relatif terhadap basis ini,
\begin{align*}
  \vec{v} &=v_1\vec{\rho}_1+\dots+v_{i-1}\vec{\rho}_{i-1}+v_i\vec{\beta}
            +v_{i+1}\vec{\rho}_{i+1}+\dots+v_n\vec{\rho}_n                \\
  \vec{w} &= w_1\vec{\rho}_1+\dots+w_{i-1}\vec{\rho}_{i-1}+w_i\vec{\beta}
            +w_{i+1}\vec{\rho}_{i+1}+\dots+w_n\vec{\rho}_n
\end{align*}
lalu jumlahkan.
\begin{equation*}
  \vec{v}+\vec{w}
  =
  (v_1+w_1)\vec{\rho}_1+\dots+(v_i+w_i)\vec{\beta}
            +\dots+(v_n+w_n)\vec{\rho}_n 
\end{equation*}
Tinjaulah ruas kiri (1), lalu uraikan $\vec{v}+\vec{w}$.
% $\det (\vec{\rho}_1,\dots,\vec{v}+\vec{w},\dots,\vec{\rho}_n)$.
\begin{equation*}
    \det (\vec{\rho}_1,\dots,\,(v_1+w_1)\vec{\rho}_1+\dots+(v_i+w_i)\vec{\beta}
            +\dots+(v_n+w_n)\vec{\rho}_n,\,\dots,\vec{\rho}_n) 
   \tag{$*$}
\end{equation*}
% Gauss's Method enables us to eliminate the terms for $\vec{\rho}_1$,
% $\vec{\rho}_2$, etc., from the expanded expression.
% the $(v_1+w_1)\vec{\rho}_1+\dots+(v_i+w_i)\vec{\beta}
%             +\dots+(v_n+w_n)\vec{\rho}_n$ expression.
Menurut syarat~(1) dalam definisi determinan, nilai 
($*$) % the determinant
% $\det(\vec{\rho}_1,\dots,\vec{v}+\vec{w},\dots,\vec{\rho}_n)$
tidak berubah oleh operasi 
penambahan $-(v_1+w_1)\vec{\rho}_1$ ke baris ke-$i$, $\vec{v}+\vec{w}$.
Baris ke-$i$ menjadi sebagai berikut.
\begin{equation*}
  \vec{v}+\vec{w}-(v_1+w_1)\vec{\rho}_1
  =
  (v_2+w_2)\vec{\rho}_2+\cdots+(v_i+w_i)\vec{\beta}
            +\dots+(v_n+w_n)\vec{\rho}_n 
\end{equation*}
%</pf:DetsMultilinear2>
%<*pf:DetsMultilinear3>
Selanjutnya tambahkan $-(v_2+w_2)\vec{\rho}_2$, dan seterusnya, untuk mengeliminasi 
semua suku dari baris-baris lainnya.
Terapkan syarat~(3) dari definisi determinan.
\begin{multline*}
  \det (\vec{\rho}_1,\dots,\vec{v}+\vec{w},\dots,\vec{\rho}_n)          \\
  \begin{aligned}
    &=\det (\vec{\rho}_1,\dots,(v_i+w_i)\cdot\vec{\beta},\dots,\vec{\rho}_n) \\
    &=(v_i+w_i)\cdot\det (\vec{\rho}_1,\dots,\vec{\beta},\dots,\vec{\rho}_n) \\
    &=v_i\cdot \det (\vec{\rho}_1,\dots,\vec{\beta},\dots,\vec{\rho}_n)  
     +w_i\cdot \det (\vec{\rho}_1,\dots,\vec{\beta},\dots,\vec{\rho}_n)
  \end{aligned}
\end{multline*}
Sekarang bentuk ini merupakan jumlah dua determinan.
Sebagai langkah terakhir, masukkan kembali $v_i$ dan $w_i$ ke depan 
$\vec{\beta}$,
lalu gunakan kembali kombinasi baris, 
kali ini untuk merekonstruksi ekspresi $\vec{v}$
dan $\vec{w}$ relatif terhadap basis tersebut.
Dengan kata lain, mulailah dengan operasi
penambahan $v_1\vec{\rho}_1$ ke $v_i\vec{\beta}$ 
dan $w_1\vec{\rho}_1$ ke $w_i\vec{\beta}$, dan seterusnya,
untuk memperoleh uraian $\vec{v}$ dan $\vec{w}$.
%</pf:DetsMultilinear3>
\end{proof}

Multilinearitas memungkinkan kita menguraikan suatu determinan menjadi jumlah beberapa
determinan, yang masing-masing melibatkan matriks sederhana.

\begin{example}
Gunakan sifat~(1) multilinearitas untuk
memecah baris pertama,
\begin{equation*}
  \begin{vmat}[r]
     2  &1  \\
     4  &3
  \end{vmat}
  =
  \begin{vmat}[r]
     2  &0  \\
     4  &3
  \end{vmat}
  +
  \begin{vmat}[r]
     0  &1  \\
     4  &3
  \end{vmat}
\end{equation*}
lalu gunakan kembali (1) untuk memecah masing-masing berdasarkan baris kedua.
\begin{equation*}
  =\begin{vmat}[r]
     2  &0  \\
     4  &0
  \end{vmat}
  +
  \begin{vmat}[r]
     2  &0  \\
     0  &3
  \end{vmat}
  +
  \begin{vmat}[r]
     0  &1  \\
     4  &0
  \end{vmat}
  +
  \begin{vmat}[r]
     0  &1  \\
     0  &3
  \end{vmat}
\end{equation*}
Hasilnya adalah empat determinan.
Pada setiap baris dari keempat determinan tersebut terdapat
satu entri dari matriks semula.
\end{example}

\begin{example}
Dengan cara yang sama,
determinan \( \nbyn{3} \) terurai menjadi jumlah 
banyak determinan yang lebih sederhana.
Pemecahan berdasarkan baris pertama menghasilkan tiga determinan
(kita menyoroti nol pada posisi $1,3$ agar secara visual dapat dibedakan
dari nol-nol yang muncul sebagai bagian dari pemecahan).
\begin{equation*}
  \begin{vmat}[r]
     2              &1  &-1  \\
     4              &3  &\highlight{0}  \\
     2              &1  &5
  \end{vmat}
  =
  \begin{vmat}[r]
     2              &0  &0   \\
     4              &3  &\highlight{0}  \\
     2              &1  &5
  \end{vmat}
  +
  \begin{vmat}[r]
     0              &1  &0   \\
     4              &3  &\highlight{0}   \\
     2              &1  &5
  \end{vmat}
  +
  \begin{vmat}[r]
     0              &0  &-1  \\
     4              &3  &\highlight{0}  \\
     2  &1  &5
  \end{vmat}                       
\end{equation*}
Selanjutnya, masing-masing determinan di atas terpecah menjadi tiga berdasarkan baris kedua.
Kemudian masing-masing dari kesembilan determinan itu terpecah menjadi tiga berdasarkan baris ketiga.
Hasilnya adalah dua puluh tujuh determinan,
sedemikian sehingga setiap baris memuat satu entri dari matriks awal.
\begin{equation*}
  =
  \begin{vmat}[r]
     2              &0  &0   \\
     4              &0  &0   \\
     2              &0  &0
  \end{vmat}
  +
  \begin{vmat}[r]
     2  &0  &0   \\
     4  &0  &0   \\
     0  &1  &0
  \end{vmat}
  +
  \begin{vmat}[r]
     2  &0  &0   \\
     4  &0  &0   \\
     0  &0  &5
  \end{vmat}
  +
  \begin{vmat}[r]
     2  &0  &0   \\
     0  &3  &0   \\
     2  &0  &0
  \end{vmat}
  +\dots+
  \begin{vmat}[r]
     0  &0  &-1  \\
     0  &0  &\highlight{0}  \\
     0  &0  &5
  \end{vmat}
\end{equation*}
\end{example}

Jadi, multilinearitas akan menguraikan determinan \( \nbyn{n} \) menjadi
jumlah sebanyak \( n^n \) determinan, dengan
setiap baris dari setiap determinan memuat satu entri dari
matriks awal.

Dalam ekspansi ini, meskipun terdapat banyak suku, 
sebagian besar determinannya bernilai nol.

\begin{example} \label{ex:SamplePermExp}
Dalam setiap contoh berikut dari ekspansi sebelumnya,
dua entri dari matriks semula berada pada kolom yang sama.
\begin{equation*}
  \begin{vmat}[r]
     2               &0  &0   \\
     4               &0  &0  \\
     0               &1  &0
  \end{vmat}
  \qquad
  \begin{vmat}[r]
     0               &0  &-1  \\
     0               &3  &0  \\
     0               &0  &5
  \end{vmat}
  \qquad\
  \begin{vmat}[r]
     0               &1  &0   \\
     0               &0  &\highlight{0}  \\
     0               &0  &5
  \end{vmat}
\end{equation*}
Sebagai contoh,
dalam matriks pertama, $2$ dan $4$ sama-sama berasal dari kolom pertama
matriks semula.
Dalam matriks kedua, $-1$ dan~$5$ sama-sama berasal dari kolom ketiga.
Dalam matriks ketiga, $0$ dan~$5$ juga sama-sama berasal dari kolom ketiga.
Setiap matriks semacam itu singular karena 
salah satu baris merupakan kelipatan baris lainnya.
Jadi, berdasarkan \nearbylemma{le:IdenRowsDetZero}, setiap determinan semacam itu bernilai nol.

Berdasarkan pengamatan tersebut, 
ekspansi determinan \( \nbyn{3} \) di atas menjadi
jumlah dua puluh tujuh determinan menyederhana menjadi jumlah enam determinan berikut,
yang entri-entri dari matriks semulanya muncul satu per baris
dan juga satu per kolom.
\begin{align*}
  \begin{vmat}[r]
     2  &1  &-1  \\
     4  &3  &\highlight{0}  \\
     2  &1  &5
  \end{vmat}
  &=\begin{vmat}[r]
     2  &0  &0   \\
     0  &3  &0   \\
     0  &0  &5
  \end{vmat}
  +
   \begin{vmat}[r]
     2  &0  &0   \\
     0  &0  &\highlight{0}   \\
     0  &1  &0
  \end{vmat}                      \\
  &\quad\hbox{}+\begin{vmat}[r]
     0  &1  &0   \\
     4  &0  &0   \\
     0  &0  &5
  \end{vmat}
  +
   \begin{vmat}[r]
     0  &1  &0   \\
     0  &0  &\highlight{0}   \\
     2  &0  &0
  \end{vmat}                      \\
  &\quad\hbox{}+\begin{vmat}[r]
     0  &0  &-1  \\
     4  &0  &0   \\
     0  &1  &0
  \end{vmat}
  +
   \begin{vmat}[r]
     0  &0  &-1  \\
     0  &3  &0    \\
     2  &0  &0
  \end{vmat}                      
\end{align*}
Dalam ekspansi tersebut kita dapat mengeluarkan faktor-faktor skalarnya.
\begin{align*}
  &=(2)(3)(5)\begin{vmat}[r]
               1  &0  &0  \\
               0  &1  &0  \\
               0  &0  &1
            \end{vmat}
  +(2)(\highlight{0})(1)\begin{vmat}[r]
               1  &0  &0  \\
               0  &0  &1  \\
               0  &1  &0
            \end{vmat}                     \\
  &\quad\hbox{}+(1)(4)(5)\begin{vmat}[r]
               0  &1  &0  \\
               1  &0  &0  \\
               0  &0  &1
            \end{vmat}
  +(1)(\highlight{0})(2)\begin{vmat}[r]
               0  &1  &0  \\
               0  &0  &1  \\
               1  &0  &0
            \end{vmat}                       \\
  &\quad\hbox{}+(-1)(4)(1)\begin{vmat}[r]
               0  &0  &1  \\
               1  &0  &0  \\
               0  &1  &0
            \end{vmat}
  +(-1)(3)(2)\begin{vmat}[r]
               0  &0  &1  \\
               0  &1  &0  \\
               1  &0  &0
            \end{vmat}                       
\end{align*}
Sebagai langkah terakhir, hitung keenam determinan tersebut dengan pertukaran baris 
hingga menjadi matriks identitas
sambil mencatat perubahan tandanya.
\begin{align*}
  &=30\cdot (+1)+0\cdot (-1)  \\
  &\quad\hbox{}+20\cdot (-1)+0\cdot (+1) \\
  &\quad \hbox{}-4\cdot (+1)-6\cdot (-1)=12
\end{align*}
\end{example}

Contoh tersebut merangkum skema perhitungan baru dalam subbagian ini.
Multilinearitas menguraikan suatu determinan menjadi 
banyak determinan terpisah, masing-masing dengan satu entri dari 
matriks semula pada setiap baris.
Sebagian besar determinan ini memiliki
satu baris yang merupakan kelipatan baris lain,
sehingga kita menghilangkannya.
Yang tersisa adalah determinan dengan satu entri per baris dan
per kolom dari matriks semula.
Mengeluarkan faktor-faktor skalar selanjutnya mereduksi
determinan yang harus kita hitung menjadi
matriks dengan satu entri per baris dan per kolom, dengan semua entri tersebut bernilai $1$.

%<*NotationForPermutationMatrices>
Ingat Definisi~Tiga.IV.\ref{df:PermutationMatrix}, bahwa
suatu \definend{matriks permutasi}\index{matriks permutasi}%
\index{matriks!permutasi}
merupakan matriks persegi dengan entri-entri $0$, kecuali
satu $1$ pada setiap baris dan kolom.
Sekarang kita memperkenalkan notasi untuk matriks permutasi.
%</NotationForPermutationMatrices>

\begin{definition}  \label{df:permutation}
%<*df:permutation>
Suatu \definend{permutasi-\( n \)}\index{permutasi}
adalah fungsi pada~$n$ bilangan bulat positif pertama,
$\map{\phi}{\set{1,\ldots,n}}{\set{1,\ldots,n}}$
yang bersifat satu-ke-satu dan pada.
%</df:permutation>
\end{definition}

Dalam suatu permutasi, setiap bilangan 
$1$, \ldots, $n$ muncul sebagai keluaran untuk tepat satu masukan.
Kita dapat menyatakan permutasi sebagai barisan 
$\phi=\sequence{\phi(1),\phi(2),\ldots,\phi(n)}$.

\begin{example} \label{ex:AllTwoThreePerms}
%<*ex:AllTwoPerms>
Permutasi-$2$ adalah fungsi-fungsi
$\map{\phi_1}{\set{1,2}}{\set{1,2}}$
yang diberikan oleh $\phi_1(1)=1$, $\phi_1(2)=2$, 
dan 
$\map{\phi_2}{\set{1,2}}{\set{1,2}}$
yang diberikan oleh $\phi_2(1)=2$, $\phi_2(2)=1$. 
Notasi barisan lebih singkat:
\( \phi_1=\sequence{1,2} \) dan \( \phi_2=\sequence{2,1} \).
%</ex:AllTwoPerms>
\end{example}

\begin{example} \label{ex:AllThreePerms}
%<*ex:AllThreePerms>
Dalam notasi barisan, permutasi-$3$ adalah
\( \phi_1=\sequence{1,2,3} \),
\( \phi_2=\sequence{1,3,2} \),
\( \phi_3=\sequence{2,1,3} \),
\( \phi_4=\sequence{2,3,1} \),
\( \phi_5=\sequence{3,1,2} \), dan
\( \phi_6=\sequence{3,2,1} \).
%</ex:AllThreePerms>
\end{example}

%<*AssociatePermutationWithPermutationMatrix>
Kita menyatakan vektor baris yang semua entrinya $0$, kecuali satu $1$ pada 
entri $j$, dengan $\iota_j$; jadi, $\iota_2$ berlebar empat 
adalah~$\rowvec{0 &1 &0 &0}$.
Sekarang notasi kita untuk matriks permutasi adalah sebagai berikut:
dengan setiap $\phi=\sequence{\phi(1),\ldots,\phi(n)}$,
pasangkan matriks yang baris-barisnya adalah $\iota_{\phi(1)}$, \ldots, $\iota_{\phi(n)}$.
%</AssociatePermutationWithPermutationMatrix>
Sebagai contoh, matriks yang dipasangkan dengan permutasi-$4$ \( \phi=\sequence{3,2,1,4} \) 
memiliki baris-baris $\iota$ yang bersesuaian.
\begin{equation*}
   P_{\phi}=
   \begin{mat}
      \iota_{3} \\
      \iota_{2} \\
      \iota_{1} \\
      \iota_{4} 
   \end{mat}=
  \begin{mat}[r]
     0  &0  &1  &0  \\
     0  &1  &0  &0  \\
     1  &0  &0  &0  \\
     0  &0  &0  &1
  \end{mat}
\end{equation*}

\begin{example} \label{ex:TwoThreePermsAndMatrices}
Berikut adalah matriks permutasi untuk permutasi-$2$ 
yang dicantumkan dalam \nearbyexample{ex:AllTwoThreePerms}.
\begin{equation*}
  P_{\phi_1}
  =\begin{mat}
      \iota_1 \\
      \iota_2 
   \end{mat}
  =\begin{mat}[r]
    1  &0         \\
    0  &1   
  \end{mat}
  \qquad
  P_{\phi_2}
   =\begin{mat}
      \iota_2 \\
      \iota_1 
   \end{mat}
   =\begin{mat}[r]
    0  &1         \\
    1  &0   
  \end{mat}
\end{equation*}
Sebagai contoh, baris pertama $P_{\phi_2}$ adalah 
$\iota_{\phi_2(1)}=\iota_2$ dan baris keduanya adalah $\iota_{\phi_2(2)}=\iota_1$.
\end{example}

\begin{example}
Tinjaulah permutasi-$3$
\( \phi_5=\sequence{3,1,2} \).
Matriks permutasi  
$P_{\phi_5}$ memiliki baris-baris $\iota_{\phi_5(1)}=\iota_3$, 
$\iota_{\phi_5(2)}=\iota_1$, dan $\iota_{\phi_5(3)}=\iota_2$.\begin{equation*}
  % P_{\phi_2}
  %  =\begin{mat}
  %     \iota_1 \\
  %     \iota_3 \\
  %     \iota_2 
  %  \end{mat}
  % =\begin{mat}[r]
  %    1      &0        &0        \\
  %    0      &0        &1        \\
  %    0      &1        &0
  % \end{mat}
  % \qquad
  P_{\phi_5}
   % =\begin{mat}
   %    \iota_3 \\
   %    \iota_1 \\
   %    \iota_2 
   % \end{mat}
  =\begin{mat}[r]
     0      &0        &1        \\
     1      &0        &0        \\
     0      &1        &0
  \end{mat}
\end{equation*}
\end{example}

\begin{definition} \label{df:PermutationExpansion}
%<*df:PermutationExpansion>
\definend{Ekspansi permutasi}\index{determinan!ekspansi permutasi}%
\index{ekspansi permutasi}
untuk determinan adalah
\begin{equation*}
   \begin{vmat}
      t_{1,1}  &t_{1,2}  &\ldots  &t_{1,n}  \\
      t_{2,1}  &t_{2,2}  &\ldots  &t_{2,n}  \\
               &\vdots                      \\
      t_{n,1}  &t_{n,2}  &\ldots  &t_{n,n}
   \end{vmat}
   =
   \begin{array}[t]{l}
      t_{1,\phi_1(1)}t_{2,\phi_1(2)}\cdots
           t_{n,\phi_1(n)}\deter{P_{\phi_1}}       \\[.5ex]
      \>\hbox{}+t_{1,\phi_2(1)}t_{2,\phi_2(2)}\cdots
           t_{n,\phi_2(n)}\deter{P_{\phi_2}}       \\
      \>\vdotswithin{+}                              \\
      \>\hbox{}+t_{1,\phi_k(1)}t_{2,\phi_k(2)}\cdots
           t_{n,\phi_k(n)}\deter{P_{\phi_k}} 
   \end{array}
\end{equation*}
dengan \( \phi_1,\ldots,\phi_k \) semua permutasi-\( n \).
%</df:PermutationExpansion>
\end{definition}
%<*SummationForPermutationExpansion>
Kita dapat menyatakan kembali rumus tersebut dalam 
\definend{notasi 
penjumlahan}\index{notasi penjumlahan, untuk ekspansi permutasi}
\begin{equation*}
  \deter{T}=\!\!
  \sum_{\text{permutasi\ }\phi}\!\!\!\!
     t_{1,\phi(1)}t_{2,\phi(2)}\cdots t_{n,\phi(n)}
                                 \deter{P_{\phi}}
\end{equation*}
yang dibaca,
``jumlah, atas semua permutasi \( \phi \), dari suku-suku berbentuk
\( t_{1,\phi(1)}t_{2,\phi(2)}\cdots t_{n,\phi(n)} \deter{P_{\phi}} \).''
%</SummationForPermutationExpansion>


\begin{example}
Rumus determinan $\nbyn{2}$ yang sudah dikenal diperoleh dari rumus di atas,
\begin{align*}
  \begin{vmat}
    t_{1,1}  &t_{1,2} \\
    t_{2,1}  &t_{2,2}
  \end{vmat}
  &=
  t_{1,1}t_{2,2}\cdot\deter{P_{\phi_1}}
  +
  t_{1,2}t_{2,1}\cdot\deter{P_{\phi_2}}      \\[-1.5ex]     
  &=
  t_{1,1}t_{2,2}\cdot\begin{vmat}[r]
           1  &0 \\
           0  &1
         \end{vmat}
  +
  t_{1,2}t_{2,1}\cdot\begin{vmat}[r]
            0  &1 \\
            1  &0
          \end{vmat}               \\
  &=t_{1,1}t_{2,2}-t_{1,2}t_{2,1}
\end{align*}
demikian pula rumus $\nbyn{3}$.
\begin{align*}
  \begin{vmat}
    t_{1,1}  &t_{1,2}  &t_{1,3} \\
    t_{2,1}  &t_{2,2}  &t_{2,3} \\
    t_{3,1}  &t_{3,2}  &t_{3,3} 
  \end{vmat}
  &=
  \begin{aligned}[t]
    &t_{1,1}t_{2,2}t_{3,3}\deter{P_{\phi_1}}
     +t_{1,1}t_{2,3}t_{3,2}\deter{P_{\phi_2}}
     +t_{1,2}t_{2,1}t_{3,3}\deter{P_{\phi_3}} \\
    &\hspace*{0em}
     +t_{1,2}t_{2,3}t_{3,1}\deter{P_{\phi_4}}
     +t_{1,3}t_{2,1}t_{3,2}\deter{P_{\phi_5}}
     +t_{1,3}t_{2,2}t_{3,1}\deter{P_{\phi_6}}
  \end{aligned}                                      \\
  &=
  \begin{aligned}[t]
    &t_{1,1}t_{2,2}t_{3,3}
     -t_{1,1}t_{2,3}t_{3,2}
     -t_{1,2}t_{2,1}t_{3,3}  \\
    &\hspace*{0em}
     +t_{1,2}t_{2,3}t_{3,1}
     +t_{1,3}t_{2,1}t_{3,2}
     -t_{1,3}t_{2,2}t_{3,1}
  \end{aligned}
\end{align*}
\end{example}

Menghitung determinan dengan ekspansi permutasi biasanya memerlukan 
waktu lebih lama daripada menggunakan Metode Gauss.
Namun, kita akan menggunakannya untuk membuktikan bahwa 
fungsi determinan ada.
Buktinya panjang, sehingga di sini kita hanya akan menyatakan hasilnya dan 
menunda buktinya hingga subbagian berikutnya.

\begin{theorem} \label{th:DetsExist}
\index{determinan!keberadaan} 
%<*th:DetsExist>
Untuk setiap $n$, terdapat fungsi determinan $\nbyn{n}$.
%</th:DetsExist>
\end{theorem}

Subbagian berikutnya juga memuat bukti hasil
berikut (keduanya disatukan karena kedua bukti saling bertumpang-tindih).

\begin{theorem}  \label{th:DeterminantOfAMatrixEqualsDeterminantOfTranspose}
\index{transpos!determinan}
%<*th:DeterminantOfAMatrixEqualsDeterminantOfTranspose>
Determinan suatu matriks sama dengan determinan transposnya.
%</th:DeterminantOfAMatrixEqualsDeterminantOfTranspose>
\end{theorem}

Berdasarkan teorema ini,
meskipun sejauh ini kita telah menyatakan hasil-hasil determinan dalam baris,
semua hasil tersebut juga berlaku dalam kolom.

\begin{corollary} \label{cor:ColSwapChgSign} \label{cor:DetsMultiInCols}
%<*co:ColSwapChgSign>
  Matriks dengan dua kolom yang sama bersifat singular.
  Pertukaran kolom mengubah tanda determinan.
  Determinan bersifat multilinear pada kolom-kolomnya.
%</co:ColSwapChgSign>
\end{corollary}

\begin{proof}
%<*pf:ColSwapChgSign>
Untuk pernyataan pertama, 
mentransposkan matriks menghasilkan matriks dengan determinan yang sama
dan dengan dua baris yang sama, sehingga determinannya nol.
Buktikan dua pernyataan lainnya dengan cara yang sama.
%</pf:ColSwapChgSign>
\end{proof}

Kita menutup subbagian ini dengan suatu ringkasan:
fungsi determinan ada, bersifat tunggal, dan kita mengetahui cara menghitungnya.
Mengenai makna determinan, mungkin bait-bait berikut
\cite{Kemp} dapat membantu kita mengingatnya.
\begin{verse} \small
Determinan nol,         \\*
\ Solusi: banyak atau nihil.   \\*
Determinan tak nol,         \\*
\ Solusi: satu saja.
\end{verse}




\begin{exercises}
  \item[]\textit{Berikut merangkum notasi kita  
    untuk permutasi-$2$\hbox{} dan permutasi-$3$.
    \begin{center}
      \begin{tabular}[t]{c|cc}
        $i$          &$1$      &$2$    \\
        \hline
        $\phi_1(i)$  &$1$      &$2$     \\
        $\phi_2(i)$  &$2$      &$1$     
      \end{tabular}
      \qquad
      \begin{tabular}[t]{c|ccc}
        $i$          &$1$     &$2$   &$3$    \\
        \hline
        $\phi_1(i)$  &$1$     &$2$   &$3$    \\
        $\phi_2(i)$  &$1$     &$3$   &$2$    \\
        $\phi_3(i)$  &$2$     &$1$   &$3$    \\
        $\phi_4(i)$  &$2$     &$3$   &$1$    \\
        $\phi_5(i)$  &$3$     &$1$   &$2$    \\
        $\phi_6(i)$  &$3$     &$2$   &$1$    
      \end{tabular}
    \end{center}  }
  \recommended \item Untuk matriks berikut, tentukan suku yang berkaitan dengan 
    setiap permutasi-$3$.
    \begin{equation*}
      M=\begin{mat}
        1 &2 &3 \\
        4 &5 &6 \\
        7 &8 &9
      \end{mat}
    \end{equation*}
    Dengan kata lain, lengkapi tabel berikut.
      \begin{center} \small
        \begin{tabular}{r|cccccc}
          \textit{permutasi} $\phi_i$ 
            &$\phi_1$ &$\phi_2$ &$\phi_3$ &$\phi_4$ &$\phi_5$ &$\phi_6$  \\
          \hline
          \textit{suku} $m_{1,\phi_i(1)}m_{2,\phi_i(2)}m_{3,\phi_i(3)}$
           &$1\cdot 5\cdot 9$ 
           & 
           & 
           & 
           & 
           & 
        \end{tabular}
      \end{center}
    \begin{answer}
      Sebut matriks tersebut~$M$.
      \begin{center}
        \begin{tabular}{r|cccccc}
          \textit{permutasi} 
            &$\phi_1$ &$\phi_2$ &$\phi_3$ &$\phi_4$ &$\phi_5$ &$\phi_6$  \\
          \hline
          \textit{suku} 
           &$1\cdot 5\cdot 9$ 
           &$1\cdot 6\cdot 8$ 
           &$2\cdot 4\cdot 9$ 
           &$2\cdot 6\cdot 7$ 
           &$3\cdot 4\cdot 8$ 
           &$3\cdot 5\cdot 7$ 
        \end{tabular}
      \end{center}
    \end{answer}
  \recommended \item Untuk setiap permutasi-$3$~$\phi$, tentukan $|P_\phi|$.
    \begin{answer}
      Kita dapat mempertukarkan baris setiap $P_{\phi}$ hingga menjadi matriks identitas.
      Jika banyaknya pertukaran genap, determinannya adalah~$+1$,
      sedangkan jika banyaknya pertukaran ganjil, determinannya adalah~$-1$.
      \begin{center}
        \begin{tabular}{r|cccccc}
          \textit{permutasi} 
            &$\phi_1$ &$\phi_2$ &$\phi_3$ &$\phi_4$ &$\phi_5$ &$\phi_6$  \\
          \hline
          \textit{banyaknya pertukaran} &$0$
                          &$1$
                          &$1$
                          &$2$
                          &$2$
                          &$1$  \\
          $|P_{\phi_i}|$      
           &$+1$ 
           &$-1$ 
           &$-1$ 
           &$+1$ 
           &$+1$ 
           &$-1$ 
        \end{tabular}
      \end{center}
      \textit{Catatan.}
      Dalam tabel tersebut 
      kita cukup melakukan pertukaran hingga menemukan suatu jumlah yang membawa kita kembali ke 
      `$1,2,3$'.
      Namun, mungkinkah banyaknya pertukaran berbeda?
      Jika seseorang menemukan $2$ pertukaran dan orang lain menemukan
      $4$, hal itu tidak menjadi masalah karena keduanya memberikan determinan~$+1$.
      Akan tetapi, jika kita dapat melakukan pertukaran melalui dua cara berbeda, yang satu genap dan yang
      lain ganjil, hal itu akan menjadi masalah.
      Bagian berikutnya menunjukkan bahwa campuran genap dan ganjil tidak mungkin terjadi.
    \end{answer}
  \item Determinan berikut bernilai~$7$ menurut rumus $\nbyn{2}$.
    Hitunglah dengan ekspansi permutasi.
    \begin{equation*}
      \begin{vmatrix}
        2 &3 \\
        1 &5
      \end{vmatrix}
    \end{equation*}
    \begin{answer}
      \begin{align*}
        \begin{vmat}
          2  &3 \\
          1  &5
        \end{vmat}
        &=
        2\cdot 5\cdot\deter{P_{\phi_1}}
        +
        1\cdot 3\cdot\deter{P_{\phi_2}}      \\     
        &=
        2\cdot 5\cdot\begin{vmat}[r]
                           1  &0 \\
                           0  &1
                         \end{vmat}
        +
        1\cdot 3\cdot\begin{vmat}[r]
                           0  &1 \\
                           1  &0
                         \end{vmat}               \\
        &=2\cdot 5\cdot 1+1\cdot 3\cdot (-1)=7
        \end{align*}
    \end{answer}
  \item Determinan berikut bernilai $0$ karena jumlah dua baris pertamanya sama dengan baris ketiga.
    Hitung determinan tersebut menggunakan ekspansi permutasi.
    \begin{equation*}
      \begin{vmatrix}
        -1 &0 &1 \\
        3  &1 &4 \\
        2  &1 &5
      \end{vmatrix}
    \end{equation*}
    \begin{answer}
          \begin{align*}
            \begin{vmat}[r]
              -1  &0  &1  \\
               3  &1  &4  \\
               2  &1  &5
            \end{vmat}
            &=
            \begin{aligned}[t]
              &(-1)(1)(5)\deter{P_{\phi_1}}
              +(-1)(4)(1)\deter{P_{\phi_2}}
              +(0)(3)(5)\deter{P_{\phi_3}}  \\
              &\hbox{}\quad\hbox{}
              +(0)(4)(2)\deter{P_{\phi_4}}
              +(1)(3)(1)\deter{P_{\phi_5}}
              +(1)(1)(2)\deter{P_{\phi_6}}
            \end{aligned}                        \\
            &=
            \begin{aligned}[t]
              &(-1)(1)(5)\cdot 1
              +(-1)(4)(1)\cdot (-1)
              +(0)(3)(5)\cdot (-1)  \\
              &\hbox{}\quad\hbox{}
              +(0)(4)(2)\cdot 1
              +(1)(3)(1)\cdot 1
              +(1)(1)(2)\cdot (-1)
            \end{aligned}                        \\
            &=-5+4+0+0+3-2=0
          \end{align*}      
    \end{answer}
  \recommended \item 
    Hitung determinan dengan menggunakan ekspansi permutasi.
    \begin{exparts*}
      \partsitem 
        $\begin{vmat}[r]
          1  &2  &3  \\
          4  &5  &6  \\
          7  &8  &9
        \end{vmat}$
      \partsitem 
        $\begin{vmat}[r]
          2  &2  &1  \\
          3  &-1 &0  \\
          -2 &0  &5
        \end{vmat}$
    \end{exparts*}
    \begin{answer}
      \begin{exparts}
        \partsitem Matriks ini singular. 
          \begin{align*}
            \begin{vmat}[r]
              1  &2  &3  \\
              4  &5  &6  \\
              7  &8  &9
            \end{vmat}
            &=
            \begin{aligned}[t]
              &(1)(5)(9)\deter{P_{\phi_1}}
              +(1)(6)(8)\deter{P_{\phi_2}}
              +(2)(4)(9)\deter{P_{\phi_3}}  \\
              &\hbox{}\quad\hbox{}
              +(2)(6)(7)\deter{P_{\phi_4}}
              +(3)(4)(8)\deter{P_{\phi_5}}
              +(7)(5)(3)\deter{P_{\phi_6}}
            \end{aligned}                        \\
            &=0
          \end{align*}
        \partsitem Matriks ini nonsingular. 
          \begin{align*}
            \begin{vmat}[r]
              2  &2  &1  \\
              3  &-1 &0  \\
              -2 &0  &5
            \end{vmat}
            &=
            \begin{aligned}[t]
              &(2)(-1)(5)\deter{P_{\phi_1}}
              +(2)(0)(0)\deter{P_{\phi_2}}
              +(2)(3)(5)\deter{P_{\phi_3}} \\
              &\hbox{}\quad\hbox{} 
              +(2)(0)(-2)\deter{P_{\phi_4}}
              +(1)(3)(0)\deter{P_{\phi_5}}
              +(-2)(-1)(1)\deter{P_{\phi_6}}
            \end{aligned}                      \\
            &=-42
        \end{align*}
      \end{exparts}
    \end{answer}
  \recommended \item
    Hitung keduanya dengan Metode Gauss dan juga dengan rumus
    ekspansi permutasi.
    \begin{exparts*}
      \partsitem \( \begin{vmat}[r]
                 2  &1  \\
                 3  &1
               \end{vmat}   \)
      \partsitem \( \begin{vmat}[r]
                  0  &1  &4  \\
                  0  &2  &3  \\
                  1  &5  &1
               \end{vmat}  \)
    \end{exparts*}
    \begin{answer}
      \begin{exparts}
        \partsitem Metode Gauss memberikan hasil berikut.
          \begin{equation*}
             \begin{vmat}[r]
                 2  &1  \\
                 3  &1
             \end{vmat}
             =
             \begin{vmat}[r]
                 2  &1  \\
                 0  &-1/2
             \end{vmat}
             =-1
          \end{equation*}
          Ekspansi permutasi memberikan hasil berikut.
          \begin{equation*}
             \begin{vmat}[r]
                 2  &1  \\
                 3  &1
             \end{vmat}
             =
             \begin{vmat}[r]
                 2  &0  \\
                 0  &1
             \end{vmat} +
             \begin{vmat}[r]
                 0  &1  \\
                 3  &0
             \end{vmat}             
             =
             (2)(1)\begin{vmat}[r]
                 1  &0  \\
                 0  &1
             \end{vmat} +
             (1)(3)\begin{vmat}[r]
                 0  &1  \\
                 1  &0
             \end{vmat}             
             =
             -1
          \end{equation*}
        \partsitem Metode Gauss memberikan hasil berikut.
          \begin{equation*}
            \begin{vmat}[r]
              0  &1  &4  \\
              0  &2  &3  \\
              1  &5  &1
            \end{vmat}
            =
            -\begin{vmat}[r]
              1  &5  &1  \\
              0  &2  &3  \\
              0  &1  &4  
            \end{vmat}
            =
            -\begin{vmat}[r]
              1  &5  &1  \\
              0  &2  &3  \\
              0  &0  &5/2  
            \end{vmat}
            =-5
          \end{equation*}
          Ekspansi permutasi memberikan hasil berikut.
          \begin{align*}
            \begin{vmat}[r]
              0  &1  &4  \\
              0  &2  &3  \\
              1  &5  &1
            \end{vmat}
            &=
            \begin{aligned}[t]
            &(0)(2)(1)\deter{P_{\phi_1}}
            +(0)(3)(5)\deter{P_{\phi_2}}
            +(1)(0)(1)\deter{P_{\phi_3}}  \\
            &\hbox{}\quad\hbox{}
            +(1)(3)(1)\deter{P_{\phi_4}}
            +(4)(0)(5)\deter{P_{\phi_5}}
            +(4)(2)(1)\deter{P_{\phi_6}}
            \end{aligned}                         \\
            &=-5
          \end{align*}
      \end{exparts}  
    \end{answer}
  \recommended \item 
    Gunakan rumus ekspansi permutasi untuk menurunkan
    rumus determinan \( \nbyn{3} \).
    \begin{answer}
        Mengikuti \nearbyexample{ex:SamplePermExp} memberikan hasil berikut.
          \begin{align*}
             \begin{vmat}
                t_{1,1}  &t_{1,2}  &t_{1,3}  \\
                t_{2,1}  &t_{2,2}  &t_{2,3}  \\
                t_{3,1}  &t_{3,2}  &t_{3,3}
             \end{vmat}
             &=\begin{aligned}[t]
                 &t_{1,1}t_{2,2}t_{3,3}\deter{P_{\phi_1}}
                   +t_{1,1}t_{2,3}t_{3,2}\deter{P_{\phi_2}} \\
                 &\hbox{}\quad\hbox{}
                  +t_{1,2}t_{2,1}t_{3,3}\deter{P_{\phi_3}}
                             +t_{1,2}t_{2,3}t_{3,1}\deter{P_{\phi_4}} \\
                 &\hbox{}\quad\hbox{}
                  +t_{1,3}t_{2,1}t_{3,2}\deter{P_{\phi_5}}
                             +t_{1,3}t_{2,2}t_{3,1}\deter{P_{\phi_6}}
               \end{aligned}                                               \\
             &=\begin{aligned}[t]
                 & t_{1,1}t_{2,2}t_{3,3}(+1)
                   +t_{1,1}t_{2,3}t_{3,2}(-1)  \\
                 &\hbox{}\quad\hbox{}
                   +t_{1,2}t_{2,1}t_{3,3}(-1)
                    +t_{1,2}t_{2,3}t_{3,1}(+1) \\
                 &\hbox{}\quad\hbox{}
                   +t_{1,3}t_{2,1}t_{3,2}(+1)
                         +t_{1,3}t_{2,2}t_{3,1}(-1)
               \end{aligned}                                               
          \end{align*}  
     \end{answer}
  \item 
    Cantumkan semua permutasi-$4$.
    \begin{answer}
      Berikut semua permutasi dengan $\phi(1)=1$,
      \begin{align*}
        &\phi_1=\sequence{1,2,3,4}
        \quad
        \phi_2=\sequence{1,2,4,3}
        \quad
        \phi_3=\sequence{1,3,2,4}  \\
        &\quad
        \phi_4=\sequence{1,3,4,2}
        \quad
        \phi_5=\sequence{1,4,2,3}
        \quad
        \phi_6=\sequence{1,4,3,2}
      \end{align*}
      yang berikutnya dengan $\phi(1)=2$,
      \begin{align*}
        &\phi_7=\sequence{2,1,3,4}
        \quad
        \phi_8=\sequence{2,1,4,3}
        \quad
        \phi_9=\sequence{2,3,1,4}    \\
        &\quad
        \phi_{10}=\sequence{2,3,4,1}
        \quad
        \phi_{11}=\sequence{2,4,1,3}
        \quad
        \phi_{12}=\sequence{2,4,3,1}
      \end{align*}
      yang berikutnya dengan $\phi(1)=3$,
      \begin{align*}
        &\phi_{13}=\sequence{3,1,2,4}
        \quad
        \phi_{14}=\sequence{3,1,4,2}
        \quad
        \phi_{15}=\sequence{3,2,1,4}  \\
        &\quad
        \phi_{16}=\sequence{3,2,4,1}
        \quad
        \phi_{17}=\sequence{3,4,1,2}
        \quad
        \phi_{18}=\sequence{3,4,2,1}
      \end{align*}
      dan yang berikutnya dengan $\phi(1)=4$.
      \begin{align*}
        &\phi_{19}=\sequence{4,1,2,3}
        \quad
        \phi_{20}=\sequence{4,1,3,2}
        \quad
        \phi_{21}=\sequence{4,2,1,3}  \\
        &\quad
        \phi_{22}=\sequence{4,2,3,1}
        \quad
        \phi_{23}=\sequence{4,3,1,2}
        \quad
        \phi_{24}=\sequence{4,3,2,1}
      \end{align*}
    \end{answer}
  \item 
    Suatu permutasi, dipandang sebagai fungsi dari himpunan
    $\set{1,..,n}$ ke dirinya sendiri, bersifat satu-ke-satu dan pada.
    Karena itu, setiap permutasi memiliki invers. 
    \begin{exparts}
      \partsitem Tentukan invers setiap permutasi-$2$.
      \partsitem Tentukan invers setiap permutasi-$3$.
    \end{exparts}
    \begin{answer}
      Masing-masing hasil berikut mudah diperiksa.
      \begin{exparts}
        \partsitem 
          \begin{tabular}[t]{r|cc}
            \textit{permutasi} &$\phi_1$  &$\phi_2$ \\
             \hline
            \textit{invers}     &$\phi_1$  &$\phi_2$ 
          \end{tabular}
        \partsitem 
          \begin{tabular}[t]{r|cccccc}
            \textit{permutasi} 
              &$\phi_1$ &$\phi_2$ &$\phi_3$ &$\phi_4$ &$\phi_5$ &$\phi_6$ \\
            \hline
            \textit{invers}
              &$\phi_1$ &$\phi_2$ &$\phi_3$ &$\phi_5$ &$\phi_4$ &$\phi_6$ 
          \end{tabular}
      \end{exparts}
    \end{answer}
  \item 
     Buktikan bahwa \( f \) bersifat multilinear jika dan hanya jika untuk semua
     \( \vec{v},\vec{w}\in V \) dan \( k_1,k_2\in\Re \), pernyataan berikut berlaku.
      \begin{equation*}
         f(\vec{\rho}_1,\dots,k_1\vec{v}_1+k_2\vec{v}_2,
           \dots,\vec{\rho}_n)
         =
         k_1f(\vec{\rho}_1,\dots,\vec{v}_1,\dots,\vec{\rho}_n)+
         k_2f(\vec{\rho}_1,\dots,\vec{v}_2,\dots,\vec{\rho}_n)
       \end{equation*}
       \begin{answer}
         Untuk arah `jika', syarat pertama dari 
         \nearbydefinition{def:multilinear} diperoleh dengan mengambil $k_1=k_2=1$,
         sedangkan syarat kedua diperoleh dengan mengambil $k_2=0$.

         Arah `hanya jika' juga rutin.
         Dari
         $
           f(\vec{\rho}_1,\dots,k_1\vec{v}_1+k_2\vec{v}_2,
             \dots,\vec{\rho}_n)
         $
         syarat pertama dari \nearbydefinition{def:multilinear} memberikan
         $
           =
           f(\vec{\rho}_1,\dots,k_1\vec{v}_1,\dots,\vec{\rho}_n)+
           f(\vec{\rho}_1,\dots,k_2\vec{v}_2,\dots,\vec{\rho}_n)
         $
         dan syarat kedua, yang diterapkan dua kali, memberikan 
         hasil yang dikehendaki. 
       \end{answer}
  \item 
    Bagaimana determinan akan berubah jika kita mengubah sifat~(4) dalam
    definisi menjadi \( \deter{I}=2 \)?
    \begin{answer}
       Semuanya akan menjadi dua kali lipat.
    \end{answer}
  \item 
    Verifikasikan pernyataan kedua dan ketiga 
    dalam \nearbycorollary{cor:ColSwapChgSign}.
    \begin{answer}
      Untuk pernyataan kedua, 
      transposkan suatu matriks, pertukarkan barisnya, lalu transposkan kembali.
      Hasilnya adalah kolom-kolom yang dipertukarkan, dan determinannya berubah dengan faktor
      \( -1 \).
      Pernyataan ketiga serupa:~untuk suatu 
      matriks, transposkan matriks itu, terapkan multilinearitas pada bagian yang sekarang menjadi
      baris, lalu transposkan kembali matriks-matriks hasilnya.
    \end{answer}
  \recommended \item
    Tunjukkan bahwa jika suatu matriks \( \nbyn{n} \) memiliki determinan tak nol,
    maka kita dapat menyatakan setiap vektor kolom
    \( \vec{v}\in\Re^n \) sebagai kombinasi linear
    kolom-kolom matriks tersebut.
    \begin{answer}
      Matriks \( \nbyn{n} \) dengan determinan tak nol memiliki rank
      \( n \), sehingga kolom-kolomnya membentuk basis bagi \( \Re^n \).  
    \end{answer}
  \item 
    \cite{Strang}
    Benar atau salah: matriks yang entri-entrinya hanya nol
    atau satu memiliki determinan yang sama dengan nol, satu, atau negatif satu.
    \begin{answer}
      Salah.
      \begin{equation*}
        \begin{vmat}[r]
          0 &1  &1  \\
          1 &0  &1 \\
          1  &1 &0
        \end{vmat}
        =2
      \end{equation*} 
     \end{answer}
  \item 
    \begin{exparts}
      \partsitem Tunjukkan bahwa terdapat $120$ suku dalam rumus ekspansi
         permutasi suatu matriks \( \nbyn{5} \).
      \partsitem 
         Berapa banyak suku yang pasti nol jika entri \( 1,2 \) bernilai nol?
    \end{exparts}
    \begin{answer}
      \begin{exparts}
        \partsitem Untuk indeks kolom entri pada baris pertama terdapat
          lima pilihan.
          Kemudian, untuk indeks kolom entri pada baris kedua terdapat
          empat pilihan.
          Dengan melanjutkan, kita memperoleh $5\cdot 4\cdot 3\cdot 2\cdot 1=120$.
          (Lihat juga pertanyaan berikutnya.)
        \partsitem Setelah memilih kolom kedua pada baris pertama, 
          kita dapat memilih entri-entri lainnya melalui \( 4\cdot 3\cdot 2\cdot 1=24 \) 
          cara.
      \end{exparts}  
    \end{answer}
  \item 
    Berapa banyak permutasi-\( n \) yang ada?
    \begin{answer}
       \( n\cdot(n-1)\cdots 2\cdot 1=n! \) 
    \end{answer}
  \item Tunjukkan bahwa invers suatu matriks permutasi adalah transposnya.
    \begin{answer}
      \cite{SchmidtSO}
      Kita akan menunjukkan bahwa $P\trans{P}=I$; argumen untuk $\trans{P}P=I$ serupa.
      Entri $i,j$ dari $P\trans{P}$ merupakan jumlah suku-suku berbentuk
      $p_{i,k}q_{k,j}$, dengan entri-entri $\trans{P}$ dinyatakan menggunakan $q$;
      yakni, $q_{k,j}=p_{j,k}$.
      Jadi, entri $i,j$ dari $P\trans{P}$ adalah jumlah
      \(
        \sum_{k=1}^n p_{i,k}p_{j,k}
      \).
      Namun, $p_{i,k}$ biasanya $0$, sehingga $P_{i,k}P_{j,k}$ biasanya $0$. 
      $P_{i,k}$ hanya bernilai tak nol ketika nilainya $1$, 
      tetapi saat itu tidak ada $i^\prime\neq i$ lain sedemikian sehingga $P_{i^\prime,k}$ 
      bernilai tak nol (baris $i$ adalah satu-satunya baris dengan $1$ pada kolom~$k$). 
      Dengan kata lain,
      \begin{equation*}
        \sum_{k=1}^n p_{i,k}p_{j,k}=
        \begin{cases}
          1  &i=j  \\
          0  &\text{jika tidak}
        \end{cases}
      \end{equation*}
      dan ini tepat merupakan rumus entri-entri matriks identitas.
    \end{answer}
  \item 
    Matriks \( A \) disebut 
    \definend{antisimetris}\index{matriks!antisimetris}%
    \index{antisimetris} jika \( \trans{A}=-A \),
    seperti matriks berikut.
    \begin{equation*}
      A=\begin{mat}[r]
          0  &3  \\
         -3  &0
        \end{mat}
    \end{equation*}
    Tunjukkan bahwa matriks antisimetris \( \nbyn{n} \) dengan determinan
    tak nol hanya ada untuk \( n \) genap.
    \begin{answer}
      Dalam 
      \( \deter{A}=\deter{\trans{A}}=\deter{-A}=(-1)^n\deter{A} \)
      eksponen $n$ harus genap. 
    \end{answer}
  \recommended \item
    Berapakah jumlah nol terkecil, beserta penempatan
    nol-nol tersebut, yang diperlukan untuk memastikan bahwa suatu matriks \( \nbyn{4} \)
    memiliki determinan nol?
    \begin{answer}
      Menunjukkan bahwa tidak ada penempatan tiga nol yang memadai bersifat rutin.
      Empat nol memang memadai; tempatkan semuanya pada 
      baris atau kolom yang sama.
    \end{answer}
  \item
    Jika kita memiliki \( n \) titik data
    \( (x_1,y_1),(x_2,y_2),\dots\,,(x_n,y_n) \) dan ingin mencari
    polinom \( p(x)=a_{n-1}x^{n-1}+a_{n-2}x^{n-2}+\dots+a_1x+a_0 \)
    yang melalui titik-titik tersebut,
    kita dapat menyubstitusikan titik-titik itu untuk memperoleh sistem linear dengan
    \( n \)~persamaan dan \( n \)~variabel tak diketahui.
    Matriks koefisien sistem tersebut adalah
    \definend{matriks Vandermonde}.\index{matriks!Vandermonde}%
    \index{matriks Vandermonde}
    Buktikan bahwa determinan transpos matriks koefisien tersebut,
    \begin{equation*}
      \begin{vmat}
         1       &1       &\ldots   &1       \\
         x_1     &x_2     &\ldots   &x_n     \\
         {x_1}^2 &{x_2}^2 &\ldots   &{x_n}^2 \\
                 &\vdots                     \\
         {x_1}^{n-1} &{x_2}^{n-1}   &\ldots   &{x_n}^{n-1}
      \end{vmat}
    \end{equation*}
    sama dengan hasil kali, atas semua indeks \( i,j\in\set{1,\dots,n} \) dengan
    \( i<j \), dari suku-suku berbentuk \( x_j-x_i \).
    (Hal ini menunjukkan bahwa 
    determinan bernilai nol, dan sistem linear tersebut tidak memiliki solusi tunggal,
    jika dan hanya jika nilai-nilai \( x_i \) dalam data tidak semuanya berbeda.)
    \begin{answer}
      Kasus $n=3$ menunjukkan apa yang perlu dilakukan.
      Operasi kombinasi baris 
      $-x_1\rho_2+\rho_3$ dan $-x_1\rho_1+\rho_2$
      memberikan hasil berikut.
      \begin{multline*}
        \begin{vmat}
          1     &1     &1     \\
          x_1   &x_2   &x_3   \\
          x_1^2 &x_2^2 &x_3^2 
        \end{vmat}
        =
        \begin{vmat}
          1     &1             &1             \\
          x_1   &x_2           &x_3           \\
          0     &(-x_1+x_2)x_2 &(-x_1+x_3)x_3 
        \end{vmat}                                           \\
        =
        \begin{vmat}
          1     &1             &1             \\
          0     &-x_1+x_2      &-x_1+x_3      \\
          0     &(-x_1+x_2)x_2 &(-x_1+x_3)x_3 
        \end{vmat}
      \end{multline*}
      Kemudian operasi kombinasi baris $x_2\rho_2+\rho_3$ memberikan
      hasil yang dikehendaki.
      \begin{equation*}
        =
        \begin{vmat}
          1     &1             &1                     \\
          0     &-x_1+x_2      &-x_1+x_3              \\
          0     &0             &(-x_1+x_3)(-x_2+x_3) 
        \end{vmat}
        =(x_2-x_1)(x_3-x_1)(x_3-x_2)
      \end{equation*}
    \end{answer}
  \item 
    Kita dapat membagi suatu matriks menjadi  
    \definend{blok-blok},\index{matriks blok}\index{matriks!blok} %
    seperti berikut,
    \begin{equation*} 
      \begin{pmat}{rr|r}
          1  &2   &0  \\
          3  &4   &0  \\  \cline{1-3}
          0  &0   &-2 
       \end{pmat}
    \end{equation*}
    yang memperlihatkan empat blok: blok persegi $\nbyn{2}$ dan $\nbyn{1}$
    di kiri atas dan kanan bawah, serta blok nol di
    kanan atas dan kiri bawah.
    Tunjukkan bahwa jika suatu matriks dapat kita partisi sebagai
    \begin{equation*}
      T=
      \begin{pmat}{c|c}
          J   &Z_2  \\  \cline{1-2}
          Z_1 &K
       \end{pmat}
    \end{equation*}
    dengan $J$ dan $K$ matriks persegi, serta $Z_1$ dan $Z_2$ seluruhnya bernilai nol,
    maka \( \deter{T}=\deter{J}\cdot\deter{K} \).
    \begin{answer}
      Misalkan \( T \) berukuran \( \nbyn{n} \),
      \( J \) berukuran \( \nbyn{p} \),
      dan \( K \) berukuran \( \nbyn{q} \).
      Terapkan rumus ekspansi permutasi,
      \begin{equation*}
        \deter{T}=
        \sum_{\text{\scriptsize permutasi }\phi}
                t_{1,\phi(1)}t_{2,\phi(2)}\dots t_{n,\phi(n)}
                \deter{P_{\phi}}
      \end{equation*}
      Karena bagian kanan atas \( T \) seluruhnya bernilai nol, jika suatu
      \( \phi \) memiliki sekurang-kurangnya satu dari \( p+1,\dots,n \) di antara
      \( p \) nomor kolom pertamanya, \( \phi(1),\dots,\phi(p) \), maka suku yang muncul
      dari \( \phi \) adalah \( 0 \)
      (misalnya, jika \( \phi(1)=n \), maka
      \( t_{1,\phi(1)}t_{2,\phi(2)}\dots t_{n,\phi(n)} \)
      bernilai \( 0 \)).
      Jadi, rumus di atas tereduksi menjadi jumlah atas semua permutasi
      dengan dua bagian:
      pertama susun ulang \( 1,\dots,p \), lalu setelah itu muncul
      suatu permutasi dari
      \( p+1,\dots,p+q \).
      Untuk melihat bahwa hal ini memberikan \( \deter{J}\cdot\deter{K} \), gunakan sifat distributif.
      \begin{multline*}
         \bigg[\sum_{\substack{\text{\scriptsize permutasi }\phi_1 \\
                              \text{\scriptsize dari } 1,\dots,p}}
               t_{1,\phi_1(1)}\cdots t_{p,\phi_1(p)} 
               \deter{P_{\phi_1}} \bigg]                                  \\    
         \cdot
         \bigg[\sum_{\substack{\text{\scriptsize permutasi }\phi_2 \\
                              \text{\scriptsize dari } p+1,\dots,p+q}}
               t_{p+1,\phi_2(p+1)}\cdots t_{p+q,\phi_2(p+q)} 
               \deter{P_{\phi_2}}                          \bigg]
      \end{multline*}  
    \end{answer}
  \item
    Buktikan bahwa untuk setiap matriks \( \nbyn{n} \) \( T \), terdapat paling banyak
    \( n \) bilangan real berbeda \( r \) sedemikian sehingga matriks \( T-rI \)
    memiliki determinan nol
    (kita akan menggunakan hasil ini dalam Bab Lima).
    \begin{answer}
      Kasus $n=3$ menunjukkan apa yang terjadi.
      \begin{equation*}
        \deter{T-rI}
        =
        \begin{vmat}
          t_{1,1}-r  &t_{1,2}   &t_{1,3}  \\ 
          t_{2,1}    &t_{2,2}-r &t_{2,3}  \\ 
          t_{3,1}    &t_{3,2}   &t_{3,3}-r  
        \end{vmat}
      \end{equation*}
      Setiap suku dalam ekspansi permutasi memiliki tiga faktor yang diambil dari 
      entri-entri matriks (misalnya, $(t_{1,1}-r)(t_{2,2}-r)(t_{3,3}-r)$
      dan $(t_{1,1}-r)(t_{2,3})(t_{3,2})$), sehingga determinannya
      dapat dinyatakan sebagai polinom dalam $r$ berderajat $3$.
      Polinom semacam itu memiliki paling banyak $3$ akar.

      Secara umum, ekspansi permutasi menunjukkan bahwa
      determinan merupakan jumlah suku-suku yang masing-masing
      memiliki \( n \) faktor, sehingga memberikan polinom berderajat $n$.
      Polinom berderajat \( n \) memiliki paling banyak \( n \) akar. 
    \end{answer}
  \puzzle \item 
    \cite{MathMag63Q307}
    Kesembilan digit positif dapat disusun menjadi larik \( \nbyn{3} \)
    melalui \( 9! \) cara.
    Tentukan jumlah determinan semua larik tersebut.
    \begin{answer}
      \answerasgiven
      Ketika dua baris suatu determinan dipertukarkan, tanda
      determinannya berubah.
      Ketika baris-baris determinan berukuran tiga-kali-tiga dipermutasikan, diperoleh \( 3 \)
      determinan positif dan \( 3 \) determinan negatif yang nilai mutlaknya sama.
      Karena itu, \( 9! \) determinan tersebut terbagi menjadi \( 9!/6 \) kelompok, yang masing-masing
      berjumlah nol.  
    \end{answer}
  \item 
    \cite{MathMag63Q237}
    Tunjukkan bahwa
    \begin{equation*}
      \begin{vmat}
        x-2  &x-3  &x-4  \\
        x+1  &x-1  &x-3  \\
        x-4  &x-7  &x-10
      \end{vmat}=0.
    \end{equation*}
    \begin{answer}
      \answerasgiven
      Ketika elemen-elemen suatu kolom dikurangkan dari elemen-elemen
      masing-masing dua kolom lainnya, elemen-elemen dalam dua kolom determinan
      yang dihasilkan bersifat proporsional, sehingga determinannya lenyap.
      Dengan kata lain,
      \begin{equation*}
        \begin{vmat}
          2  &1    &x-4  \\
          4  &2    &x-3  \\
          6  &3    &x-10
        \end{vmat}=
        \begin{vmat}
          1    &x-3  &-1   \\
          2    &x-1  &-2   \\
          3    &x-7  &-3
        \end{vmat}=
        \begin{vmat}
          x-2  &-1   &-2   \\
          x+1  &-2   &-4   \\
          x-4  &-3   &-6
        \end{vmat}=0.
      \end{equation*}  
    \end{answer}
  \puzzle \item 
    \cite{Monthly49p33}
    Misalkan \( S \) adalah jumlah elemen-elemen bilangan bulat dalam persegi ajaib orde
    tiga dan \( D \) adalah nilai persegi tersebut ketika dipandang sebagai
    determinan.
    Tunjukkan bahwa \( D/S \) merupakan bilangan bulat.
    \begin{answer}
      \answerasgiven
      Misalkan
      \begin{equation*}
        \begin{array}{ccc}
          a  &b  &c  \\
          d  &e  &f  \\
          g  &h  &i
        \end{array}
      \end{equation*}
      memiliki jumlah ajaib \( N=S/3 \).
      Maka
      \begin{align*}
        N
        &=(a+e+i)+(d+e+f)+(g+e+c)   \\
        &\quad\text{}-(a+d+g)-(c+f+i)=3e
      \end{align*}
      dan \( S=9e \).
      Karena itu, dengan menjumlahkan baris dan kolom,
      \begin{equation*}
        D=
        \begin{vmat}
          a  &b  &c  \\
          d  &e  &f  \\
          g  &h  &i
        \end{vmat}
        =\begin{vmat}
          a  &b  &c  \\
          d  &e  &f  \\
         3e  &3e &3e
        \end{vmat}             
        =\begin{vmat}
          a  &b  &3e \\
          d  &e  &3e \\
         3e  &3e &9e
        \end{vmat}             
        =\begin{vmat}
          a  &b  &e  \\
          d  &e  &e  \\
          1  &1  &1
        \end{vmat}S.
      \end{equation*}  
    \end{answer}
  \puzzle \item 
    \cite{Monthly31p355}
    Tunjukkan bahwa determinan dari \( n^2 \) elemen di sudut kiri atas
    segitiga Pascal
    \begin{equation*}
      \begin{array}{cccccc}
        1  &1  &1  &1  &.  &.  \\
        1  &2  &3  &.  &.      \\
        1  &3  &.  &.  &   &   \\
        1  &.  &.  &   &   &   \\
        .                      \\
        .
      \end{array}
    \end{equation*}
    bernilai satu.
    \begin{answer}
      \answerasgiven
      Nyatakan determinan yang dimaksud dengan \( D_n \) dan elemen \( a_{i,j} \)
      pada baris ke-\( i \) dan kolom ke-\( j \).
      Berdasarkan aturan pembentukan elemen-elemennya, kita memperoleh
      \begin{equation*}
        a_{i,j}=a_{i,j-1}+a_{i-1,j},
        \qquad a_{1,j}=a_{i,1}=1.
      \end{equation*}
      Kurangkan setiap baris \( D_n \) dari baris setelahnya, dengan memulai
      proses dari pasangan baris terakhir.
      Setelah \( n-1 \) pengurangan, kesamaan di atas menunjukkan bahwa
      elemen
      \( a_{i,j} \) digantikan oleh elemen \( a_{i,j-1} \), dan semua
      elemen pada kolom pertama, kecuali \( a_{1,1}=1 \), menjadi nol.
      Sekarang kurangkan setiap kolom dari kolom setelahnya, dengan memulai dari
      pasangan terakhir.
      Setelah proses ini, elemen \( a_{i,j-1} \) digantikan oleh
      \( a_{i-1,j-1} \), sebagaimana diperlihatkan oleh relasi di atas.
      Hasil kedua operasi tersebut adalah mengganti \( a_{i,j} \) dengan
      \( a_{i-1,j-1} \), serta mereduksi setiap elemen pada baris pertama dan
      kolom pertama menjadi nol.
      Karena itu, \( D_n=D_{n-1} \), sehingga
      \begin{equation*}
        D_n=D_{n-1}=D_{n-2}=\dots=D_2=1.
      \end{equation*}  
    \end{answer}
\end{exercises}














\subsectionoptional{Determinan Ada}
\textit{Subbagian ini memuat bukti dua hasil dari
subbagian sebelumnya.
Subbagian ini bersifat opsional.
Kita hanya akan menggunakan materi yang dikembangkan di sini dalam
subbagian Bentuk Kanonis Jordan, yang juga bersifat opsional.}

Kita ingin menunjukkan bahwa untuk setiap ukuran~$n$,
fungsi determinan pada matriks $\nbyn{n}$ terdefinisi dengan baik.
Subbagian sebelumnya mengembangkan 
rumus ekspansi permutasi.\index{determinan!ekspansi permutasi}%
\index{ekspansi permutasi}
\begin{align*}
   %\deter{T}
   %=
   \begin{vmat}
      t_{1,1}  &t_{1,2}  &\ldots  &t_{1,n}  \\
      t_{2,1}  &t_{2,2}  &\ldots  &t_{2,n}  \\
               &\vdots                      \\
      t_{n,1}  &t_{n,2}  &\ldots  &t_{n,n}
   \end{vmat}
   &=
   \begin{array}[t]{l}
      t_{1,\phi_1(1)}t_{2,\phi_1(2)}\cdots
           t_{n,\phi_1(n)}\deter{P_{\phi_1}}       \\[.75ex]
      \>\hbox{}+t_{1,\phi_2(1)}t_{2,\phi_2(2)}\cdots
           t_{n,\phi_2(n)}\deter{P_{\phi_2}}       \\[.75ex]
      \>\alignedvdots                              \\
      \>\hbox{}+t_{1,\phi_k(1)}t_{2,\phi_k(2)}\cdots
           t_{n,\phi_k(n)}\deter{P_{\phi_k}} 
   \end{array}                                                 \\[1ex]
   &=\!\!\sum_{\text{permutasi\ }\phi}\!\!\!\!
     t_{1,\phi(1)}t_{2,\phi(2)}\cdots t_{n,\phi(n)}
                                 \deter{P_{\phi}}
\end{align*}
Hal ini mereduksi masalah pembuktian bahwa
determinan terdefinisi dengan baik 
menjadi sekadar pembuktian bahwa determinan 
terdefinisi dengan baik pada himpunan matriks permutasi.

Suatu matriks permutasi dapat diubah menjadi matriks identitas melalui pertukaran baris.
Jadi, salah satu cara untuk
menghitung determinannya adalah mencatat banyaknya pertukaran.
Namun,
kita masih harus menunjukkan bahwa hasilnya terdefinisi dengan baik.
Ingat kembali kesulitannya:~determinan
\begin{equation*}
   P_{\phi}=
   \begin{mat}[r]
      0  &1  &0  &0 \\
      1  &0  &0  &0 \\
      0  &0  &1  &0 \\
      0  &0  &0  &1
   \end{mat}
\end{equation*}
dapat dihitung dengan satu pertukaran,
\begin{equation*}
   P_{\phi}
   \grstep{\rho_1\leftrightarrow\rho_2}
   \begin{mat}[r]
      1  &0  &0  &0 \\
      0  &1  &0  &0 \\
      0  &0  &1  &0 \\
      0  &0  &0  &1
   \end{mat}
\end{equation*}
atau dengan tiga pertukaran.
\begin{equation*}
   P_{\phi}
   \grstep{\rho_3\leftrightarrow\rho_1}
   % \begin{mat}[r]
   %    0  &0  &1  &0 \\
   %    1  &0  &0  &0 \\
   %    0  &1  &0  &0 \\
   %    0  &0  &0  &1
   % \end{mat}
   \repeatedgrstep{\rho_2\leftrightarrow\rho_3}
   % \begin{mat}[r]
   %    0  &0  &1  &0 \\
   %    0  &1  &0  &0 \\
   %    1  &0  &0  &0 \\
   %    0  &0  &0  &1
   % \end{mat}
   \repeatedgrstep{\rho_1\leftrightarrow\rho_3}
   \begin{mat}[r]
      1  &0  &0  &0 \\
      0  &1  &0  &0 \\
      0  &0  &1  &0 \\
      0  &0  &0  &1
   \end{mat}
\end{equation*}
Kedua reduksi tersebut
memiliki banyak pertukaran ganjil, sehingga dalam kasus ini 
kita memperoleh \( \deter{P_{\phi}}=-1 \).
Namun, jika ada cara melakukannya dengan banyak pertukaran genap,
maka determinan akan memberikan dua keluaran berbeda
dari satu masukan.
Di bawah ini, \nearbycorollary{cor:ParityInversEqParitySwaps} membuktikan 
bahwa hal tersebut tidak dapat terjadi\Dash
tidak ada matriks permutasi
yang dapat diubah melalui pertukaran baris menjadi matriks identitas melalui dua cara, satu dengan 
banyak pertukaran genap dan yang lain dengan banyak pertukaran ganjil.

% So the critical step will be a way to calculate whether the 
% number of swaps that it takes could be even or odd.

\begin{definition} \label{df:Inversion}
%<*df:Inversion>
Dalam suatu permutasi $\phi=\sequence{\ldots,k,\ldots,j,\ldots}$,
elemen-elemen dengan $k>j$ berada dalam suatu 
\definend{inversi} dari urutan alaminya.
Serupa dengan itu, dalam suatu matriks permutasi, dua baris
\begin{equation*}
  P_{\phi}=
  \begin{mat}
    \vdots          \\
    \iota_{k} \\
    \vdots          \\
    \iota_{j} \\
    \vdots
  \end{mat}
\end{equation*}
dengan \( k>j \) berada dalam suatu \definend{inversi}.\index{inversi}%
\index{permutasi!inversi}
%</df:Inversion>
\end{definition}

\begin{example}
Matriks permutasi berikut
\begin{equation*}
  \begin{mat}[r]
    1  &0  &0  &0 \\
    0  &0  &1  &0 \\
    0  &1  &0  &0 \\
    0  &0  &0  &1
  \end{mat}
  =
  \begin{mat}
    \iota_1  \\
    \iota_3  \\
    \iota_2  \\
    \iota_4
  \end{mat}
\end{equation*}
memiliki satu inversi, yaitu \( \iota_3 \) mendahului \( \iota_2 \).
\end{example}

\begin{example} \label{ex:ThreeInversions}
Di sini terdapat tiga inversi:
\begin{equation*}
  \begin{mat}[r]
    0  &0  &1  \\
    0  &1  &0  \\
    1  &0  &0
  \end{mat}
  =
  \begin{mat}
    \iota_3  \\
    \iota_2  \\
    \iota_1
  \end{mat}
\end{equation*}
\( \iota_3 \) mendahului \( \iota_1 \),
\( \iota_3 \) mendahului \( \iota_2 \), dan \( \iota_2 \) mendahului
\( \iota_1 \).
\end{example}

\begin{lemma} \label{le:SwapsChangeSgn}
%<*lm:SwapsChangeSgn>
Pertukaran baris dalam matriks permutasi mengubah banyaknya 
inversi dari genap menjadi ganjil, atau dari ganjil menjadi genap.
%</lm:SwapsChangeSgn>
\end{lemma}

\begin{proof}
%<*pf:SwapsChangeSgn0>
Tinjaulah pertukaran baris $j$ dan $k$, dengan $k>j$.

Jika kedua baris tersebut bersebelahan,
\begin{equation*}
   P_{\phi}=
   \begin{mat}
     \vdots           \\
     \iota_{\phi(j)}  \\
     \iota_{\phi(k)}  \\
     \vdots
   \end{mat}
 \grstep{\rho_k\leftrightarrow\rho_j}
   \begin{mat}
     \vdots           \\
     \iota_{\phi(k)}  \\
     \iota_{\phi(j)}  \\
     \vdots
   \end{mat}
\end{equation*}
maka, karena inversi yang melibatkan baris di luar pasangan ini
tidak terpengaruh,
pertukaran tersebut mengubah jumlah inversi sebanyak satu, 
dengan menghapus atau menghasilkan satu inversi bergantung pada apakah
\( \phi(j)>\phi(k) \) atau tidak.
Akibatnya, jumlah inversi berubah dari ganjil menjadi genap 
atau dari genap menjadi ganjil.
%</pf:SwapsChangeSgn0>

%<*pf:SwapsChangeSgn1>
Jika kedua baris tidak bersebelahan, kita dapat mempertukarkannya 
melalui rangkaian pertukaran baris bersebelahan, mula-mula dengan membawa baris~$k$ ke atas,
\begin{equation*}
   \begin{mat}
     \vdots             \\
     \iota_{\phi(j)}    \\
     \iota_{\phi(j+1)}  \\
     \iota_{\phi(j+2)}  \\
     \vdots             \\
     \iota_{\phi(k)}    \\
     \vdots
   \end{mat}
  \grstep{\rho_k\swap\rho_{k-1}}
  \repeatedgrstep{\rho_{k-1}\swap\rho_{k-2}}
  \cdots
  \grstep{\rho_{j+1}\swap\rho_j}
   \begin{mat}
     \vdots             \\
     \iota_{\phi(k)}    \\
     \iota_{\phi(j)}    \\
     \iota_{\phi(j+1)}  \\
     \vdots             \\
     \iota_{\phi(k-1)}  \\
     \vdots
   \end{mat}
\end{equation*}
%</pf:SwapsChangeSgn1>
%<*pf:SwapsChangeSgn2>
lalu membawa baris~$j$ ke bawah.
\begin{equation*}
  \grstep{\rho_{j+1}\swap\rho_{j+2}}
  \repeatedgrstep{\rho_{j+2}\swap\rho_{j+3}}
  \cdots
  \grstep{\rho_{k-1}\swap\rho_k}
   \begin{mat}
     \vdots             \\
     \iota_{\phi(k)}    \\
     \iota_{\phi(j+1)}  \\
     \iota_{\phi(j+2)}  \\
     \vdots             \\
     \iota_{\phi(j)}    \\
     \vdots
   \end{mat}
\end{equation*}
Setiap pertukaran bersebelahan ini mengubah banyaknya 
inversi dari ganjil menjadi genap atau dari genap menjadi ganjil.
Jumlah pertukaran \( (k-j)+(k-j-1) \) bersifat ganjil.
Jadi, secara keseluruhan, banyaknya inversi berubah 
dari genap menjadi ganjil, atau dari ganjil menjadi genap.
%</pf:SwapsChangeSgn2>
\end{proof}

\begin{corollary}  \label{cor:ParityInversEqParitySwaps}
%<*co:ParityInversEqParitySwaps>
Jika suatu matriks permutasi memiliki banyak inversi ganjil, maka mengubahnya
menjadi matriks identitas memerlukan banyak pertukaran ganjil.
Jika matriks tersebut memiliki banyak inversi genap, maka mengubahnya menjadi
matriks identitas memerlukan banyak pertukaran genap.
%</co:ParityInversEqParitySwaps>
\end{corollary}

\begin{proof}
%<*pf:ParityInversEqParitySwaps>
Matriks identitas memiliki nol inversi.
Mengubah bilangan ganjil menjadi nol memerlukan banyak pertukaran ganjil,
dan mengubah bilangan genap menjadi nol memerlukan banyak pertukaran genap.
%</pf:ParityInversEqParitySwaps>
\end{proof}

\begin{example}
Matriks dalam \nearbyexample{ex:ThreeInversions}
dapat diubah menjadi matriks identitas dengan satu pertukaran $\rho_1\leftrightarrow\rho_3$.
(Jadi, banyaknya pertukaran tidak harus sama dengan banyaknya inversi, tetapi
paritas kedua bilangan tersebut sama.)
\end{example}

\begin{definition} \label{df:Signum}
%<*df:Signum>
\definend{Signum}\index{permutasi!signum}\index{signum}%
\index{sgn!lihat {signum}}
suatu permutasi \( \sgn(\phi) \)
adalah \( -1 \) jika banyaknya inversi dalam \( \phi \) ganjil, dan
\( +1 \) jika banyaknya inversi genap.
%</df:Signum>
\end{definition}

\begin{example}
Dengan menggunakan notasi permutasi-$3$
dari \nearbyexample{ex:AllTwoThreePerms}, kita memperoleh
\begin{equation*}
  P_{\phi_1}=
  \begin{mat}
    1  &0  &0  \\
    0  &1  &0  \\
    0  &0  &1
  \end{mat}
  \qquad
  P_{\phi_2}=
  \begin{mat}
    1  &0  &0  \\
    0  &0  &1  \\
    0  &1  &0  
  \end{mat}
\end{equation*}
sehingga
\( \sgn(\phi_1)=1 \) karena tidak terdapat inversi, 
sedangkan \( \sgn(\phi_2)=-1 \) karena terdapat satu inversi.
\end{example}

Kita masih belum menunjukkan bahwa fungsi determinan
terdefinisi dengan baik karena kita belum meninjau operasi baris 
pada matriks permutasi selain pertukaran baris.
Kita akan menangani persoalan ini secara tidak langsung.
%<*DefiningDFunction>
Definisikan suatu fungsi
\( \map{d}{\matspace_{\nbyn{n}}}{\Re} \)
dengan mengubah rumus ekspansi permutasi, yaitu mengganti
$\deter{P_\phi}$ dengan $\sgn(\phi)$.
\begin{equation*}
  d(T)=
  \sum_{\text{permutasi\ }\phi}t_{1,\phi(1)}t_{2,\phi(2)}\cdots t_{n,\phi(n)}
                                 \cdot\sgn(\phi)
\end{equation*}
% This gives the same value as the permutation expansion
% because the corollary shows that $\det(P_\phi)=\sgn(\phi)$.
Keunggulan rumus ini ialah bahwa banyaknya inversi 
jelas terdefinisi dengan baik\Dash cukup hitung inversinya.
Karena itu, pembuktian 
bahwa fungsi determinan~$\nbyn{n}$ ada 
akan selesai setelah kita menunjukkan bahwa~$d$
memenuhi syarat-syarat dalam definisi determinan.
%</DefiningDFunction>

\begin{lemma} \label{lm:DeterminantsExist}
\index{determinan!keberadaan}
%<*lm:DeterminantsExist>
Fungsi \( d \) di atas merupakan determinan.
Dengan demikian, fungsi determinan $\det_{\nbyn{n}}$ ada untuk setiap $n$.
%</lm:DeterminantsExist>
\end{lemma}

\begin{proof}
%<*pf:DeterminantsExist0>
Kita harus memeriksa bahwa fungsi tersebut memenuhi empat syarat
dari definisi determinan, \nearbydefinition{def:Det}.
%</pf:DeterminantsExist0>

%<*pf:DeterminantsExist1>
Syarat~(4) mudah: jika $I$ adalah matriks identitas $\nbyn{n}$, maka dalam
\begin{equation*}
  d(I)=
  \sum_{\text{permutasi\ }\phi}
    \iota_{1,\phi(1)}\iota_{2,\phi(2)}\cdots \iota_{n,\phi(n)}
                                 \sgn(\phi)
\end{equation*}
semua suku dalam penjumlahan bernilai nol,
kecuali suku ketika permutasi~$\phi$ merupakan 
identitas, yang memberikan hasil kali entri-entri diagonal,
yaitu satu.
%</pf:DeterminantsExist1>

%<*pf:DeterminantsExist2>
Untuk syarat~(3), misalkan \( T\smash[b]{\grstep{k\rho_i}}\hat{T} \), lalu
tinjaulah $d(\hat{T})$.
\begin{multline*}
  \sum_{\text{permutasi\ }\phi}\!\!
    \hat{t}_{1,\phi(1)}
     \cdots\hat{t}_{i,\phi(i)}\cdots\hat{t}_{n,\phi(n)}
                                 \sgn(\phi)   \\
  =\sum_{\phi}
     t_{1,\phi(1)}\cdots kt_{i,\phi(i)}\cdots t_{n,\phi(n)}
                                 \sgn(\phi)
\end{multline*}
Keluarkan faktor~\( k \) untuk memperoleh kesamaan yang dikehendaki.
\begin{equation*}
  =k\cdot\sum_{\phi}
     t_{1,\phi(1)}\cdots t_{i,\phi(i)}\cdots t_{n,\phi(n)}
                                 \sgn(\phi)                 
  =k\cdot d(T)
\end{equation*}
%</pf:DeterminantsExist2>

%<*pf:DeterminantsExist3>
Untuk~(2), misalkan 
\( T\smash[b]{\grstep{\rho_i\swap\rho_j}}\hat{T} \).
Kita harus menunjukkan bahwa $d(\hat{T})$ merupakan negatif dari $d(T)$.
\begin{equation*}
  d(\hat{T})=
  \sum_{\text{permutasi\ }\phi}\!\!
    \hat{t}_{1,\phi(1)}
    \cdots\hat{t}_{i,\phi(i)}
    \cdots\hat{t}_{j,\phi(j)}
    \cdots \hat{t}_{n,\phi(n)}
                                 \sgn(\phi)
    \tag{$*$}
\end{equation*}
Kita akan menunjukkan bahwa setiap suku dalam ($*$) dipasangkan dengan suatu suku dalam $d(T)$, 
dan bahwa kedua suku tersebut saling berlawanan tanda.
Tinjaulah matriks dari ekspansi multilinear 
$d(\hat{T})$ yang memberikan suku  
    $\hat{t}_{1,\phi(1)}
    \cdots\hat{t}_{i,\phi(i)}
    \cdots\hat{t}_{j,\phi(j)}
    \cdots \hat{t}_{n,\phi(n)}
                                 \sgn(\phi)$.
\begin{equation*}
  \begin{mat}
     \ &                  &\vdots &               &     \\
       &\hat{t}_{i,\phi(i)}  &       &               &    \\
       &                &\vdots &                 &  \\
       &                &       &\hat{t}_{j,\phi(j)} &  \\
       &                &\vdots                   &
  \end{mat}
\end{equation*}
%</pf:DeterminantsExist3>
%<*pf:DeterminantsExist4>
Matriks itu merupakan hasil operasi $\rho_i\swap\rho_j$ yang dilakukan pada
matriks berikut.
\begin{equation*}
  \begin{mat}
       \ &              &\vdots &           &        \\
         &              &       &t_{i,\phi(j)} &  \\
         &              &\vdots &           &       \\
         &t_{j,\phi(i)}    &       &          &        \\
         &              &\vdots &          &
  \end{mat}
\end{equation*}
Dengan kata lain, suku dengan $t$ bertopi dipasangkan dengan 
suku berikut dari ekspansi $d(T)$:
 \(    t_{1,\sigma(1)}
       \cdots t_{j,\sigma(j)}
       \cdots t_{i,\sigma(i)}
       \cdots t_{n,\sigma(n)} \sgn(\sigma)
       \), dengan permutasi
\( \sigma \) 
sama dengan \( \phi \), tetapi bilangan
ke-\( i \) dan ke-\( j \)-nya dipertukarkan,
\( \sigma(i)=\phi(j) \) dan \( \sigma(j)=\phi(i) \).
Kedua suku tersebut memiliki faktor-faktor yang sama,
$\hat{t}_{1,\phi(1)}=t_{1,\sigma(1)}$, \ldots, 
termasuk entri-entri dari baris yang dipertukarkan,
$\hat{t}_{i,\phi(i)}=t_{j,\phi(i)}=t_{j,\sigma(j)}$
dan $\hat{t}_{j,\phi(j)}=t_{i,\phi(j)}=t_{i,\sigma(i)}$. 
Namun, kedua suku tersebut saling berlawanan tanda
karena \( \sgn(\phi)=-\sgn(\sigma) \), berdasarkan \nearbylemma{le:SwapsChangeSgn}.
%</pf:DeterminantsExist4>

%<*pf:DeterminantsExist5>
Sekarang, setiap permutasi $\phi$ dapat diperoleh dari suatu 
permutasi lain $\sigma$ melalui pertukaran semacam itu,
dengan tepat satu cara.
Karena itu, penjumlahan dalam ($*$) sebenarnya merupakan jumlah atas semua permutasi,
yang masing-masing diambil tepat satu kali.
\begin{align*}
  d(\hat{T})
  &=
  \sum_{\text{permutasi\ }\phi}\!\!
    \hat{t}_{1,\phi(1)}
    \cdots\hat{t}_{i,\phi(i)}
    \cdots\hat{t}_{j,\phi(j)}
    \cdots \hat{t}_{n,\phi(n)}
                                 \sgn(\phi)   \\
  &=\sum_{\text{permutasi\ }\sigma}\!\!
     t_{1,\sigma(1)}
    \cdots t_{j,\sigma(j)}
    \cdots t_{i,\sigma(i)}
    \cdots t_{n,\sigma(n)}
                                 \cdot\bigl(-\sgn(\sigma)\bigr)        
\end{align*}
Jadi, \( d(\hat{T})=-d(T) \).
%</pf:DeterminantsExist5>

%<*pf:DeterminantsExist6>
Terakhir, untuk syarat~(1), misalkan
\( T\smash[b]{\grstep{k\rho_i+\rho_j}}\hat{T} \).
\begin{align*}
  d(\hat{T})
  &=\sum_{\text{permutasi\ }\phi}
    \hat{t}_{1,\phi(1)}\cdots\hat{t}_{i,\phi(i)}
    \cdots\hat{t}_{j,\phi(j)}\cdots\hat{t}_{n,\phi(n)}
                                 \sgn(\phi)                \\
  &=\sum_{\phi}
    t_{1,\phi(1)}\cdots t_{i,\phi(i)}
    \cdots (kt_{i,\phi(j)}+t_{j,\phi(j)})\cdots t_{n,\phi(n)}
                                 \sgn(\phi)
\end{align*}
% (notice:~that's 
% \( kt_{i,\phi(j)} \), not \( kt_{j,\phi(j)} \)).
Gunakan sifat distributif pada penjumlahan dalam $kt_{i,\phi(j)}+t_{j,\phi(j)}$.
\begin{equation*}
  = 
      \begin{aligned}[t]
        &\smash[b]{\sum_{\smash{\phi}}}\big[t_{1,\phi(1)}\cdots t_{i,\phi(i)}
          \cdots kt_{i,\phi(j)}\cdots t_{n,\phi(n)}
                                 \sgn(\phi)         \\ % [-1ex]
        &\hbox{}\qquad\quad\hbox{}
           +t_{1,\phi(1)}\cdots t_{i,\phi(i)}
           \cdots t_{j,\phi(j)}\cdots t_{n,\phi(n)}
                                 \sgn(\phi)\big]
      \end{aligned}                                               
\end{equation*}
Pecahlah menjadi dua penjumlahan.
\begin{equation*}
  =\begin{aligned}[t] 
         &\smash[b]{\sum_{\smash{\phi}}} \:t_{1,\phi(1)}\cdots t_{i,\phi(i)}
            \cdots kt_{i,\phi(j)}\cdots t_{n,\phi(n)}
                                 \sgn(\phi)           \\ % [-1ex]               
         &\hbox{}\qquad\quad\hbox{}
            +\smash[b]{\sum_{\phi}}\:
            t_{1,\phi(1)}\cdots t_{i,\phi(i)}
            \cdots t_{j,\phi(j)}\cdots t_{n,\phi(n)}
                                 \sgn(\phi)        
     \end{aligned}                                               
\end{equation*}
Kenali penjumlahan kedua.
\begin{equation*}
  =\begin{aligned}[t]
       &k\cdot \smash[b]{\sum_{\smash{\phi}}}\:
            t_{1,\phi(1)}\cdots t_{i,\phi(i)}
            \cdots t_{i,\phi(j)}\cdots t_{n,\phi(n)}
                                 \sgn(\phi)          \\  % [-.75ex]
       &\hbox{}\qquad\quad\hbox{}
             +d(T)          
     \end{aligned}
\end{equation*}
%</pf:DeterminantsExist6>
%<*pf:DeterminantsExist7>
Tinjaulah suku-suku
\( t_{1,\phi(1)}\cdots t_{i,\phi(i)} \cdots t_{i,\phi(j)}\cdots t_{n,\phi(n)}
   \sgn(\phi)  \).
Perhatikan subskripnya; entrinya adalah \( t_{i,\phi(j)} \), bukan \( t_{j,\phi(j)} \).
Jumlah suku-suku ini merupakan determinan suatu matriks~\( S \) 
yang sama dengan \( T \),
kecuali bahwa baris~$j$ dari $S$ merupakan salinan baris~$i$ dari $T$;  
artinya, $S$ memiliki dua baris yang sama.
Dengan cara yang sama seperti ketika membuktikan \nearbylemma{le:IdenRowsDetZero}, kita dapat melihat bahwa
$d(S)=0$: pertukaran dua baris $S$ yang sama akan mengubah 
tanda $d(S)$, tetapi karena matriks tersebut tidak berubah oleh pertukaran itu, 
nilai $d(S)$ juga harus tetap, sehingga nilainya harus nol.
%</pf:DeterminantsExist7>
\end{proof}

Sekarang kita telah membuktikan bahwa fungsi determinan ada untuk setiap ukuran~$\nbyn{n}$.
Kita juga telah mengetahui bahwa untuk setiap ukuran terdapat paling banyak satu
determinan.
Karena itu, untuk setiap ukuran terdapat
tepat satu fungsi determinan.

Kita mengakhiri subbagian ini dengan membuktikan hasil lain 
yang tersisa dari subbagian sebelumnya.

\begin{theorem} \label{th:DetMatrixEqualsDetTrans}
\index{transpos!determinan}
%<*th:DetMatrixEqualsDetTrans>
Determinan suatu matriks sama dengan determinan transposnya.
%</th:DetMatrixEqualsDetTrans>
\end{theorem}

\begin{proof}
%<*pf:DetMatrixEqualsDetTrans0>
Bukti ini paling mudah dipahami dengan mengerjakan kasus umum
$\nbyn{3}$.
Akan jelas bahwa argumen tersebut berlaku untuk kasus $\nbyn{n}$.

Bandingkan ekspansi permutasi matriks~$T$
\begin{align*}
  \begin{vmat}
    t_{1,1}  &t_{1,2}  &t_{1,3}  \\
    t_{2,1}  &t_{2,2}  &t_{2,3}  \\
    t_{3,1}  &t_{3,2}  &t_{3,3}  
  \end{vmat}
  &=
  t_{1,1}t_{2,2}t_{3,3}
  \begin{vmat}
    1 &0 &0 \\
    0 &1 &0 \\
    0 &0 &1
  \end{vmat}
  +t_{1,1}t_{2,3}t_{3,2}
  \begin{vmat}
    1 &0 &0 \\
    0 &0 &1 \\
    0 &1 &0
  \end{vmat}            \\
  &\quad+
  t_{1,2}t_{2,1}t_{3,3}
  \begin{vmat}
    0 &1 &0 \\
    1 &0 &0 \\
    0 &0 &1
  \end{vmat}
  +t_{1,2}t_{2,3}t_{3,1}
  \begin{vmat}
    0 &1 &0 \\
    0 &0 &1 \\
    1 &0 &0
  \end{vmat}           \\
  &\quad+
  t_{1,3}t_{2,1}t_{3,2}
  \begin{vmat}
    0 &0 &1 \\
    1 &0 &0 \\
    0 &1 &0
  \end{vmat}
  +t_{1,3}t_{2,2}t_{3,1}
  \begin{vmat}
    0 &0 &1 \\
    0 &1 &0 \\
    1 &0 &0
  \end{vmat}
\end{align*}
dengan ekspansi permutasi transposnya.
%</pf:DetMatrixEqualsDetTrans0>
%<*pf:DetMatrixEqualsDetTrans1>
\begin{align*}
  \begin{vmat}
    t_{1,1}  &t_{2,1}  &t_{3,1}  \\
    t_{1,2}  &t_{2,2}  &t_{3,2}  \\
    t_{1,3}  &t_{2,3}  &t_{3,3}  
  \end{vmat}
  &=
  t_{1,1}t_{2,2}t_{3,3}
  \begin{vmat}
    1 &0 &0 \\
    0 &1 &0 \\
    0 &0 &1
  \end{vmat}
  +t_{1,1}t_{3,2}t_{2,3}
  \begin{vmat}
    1 &0 &0 \\
    0 &0 &1 \\
    0 &1 &0
  \end{vmat}            \\
  &\quad+
  t_{2,1}t_{1,2}t_{3,3}
  \begin{vmat}
    0 &1 &0 \\
    1 &0 &0 \\
    0 &0 &1
  \end{vmat}
  +t_{2,1}t_{3,2}t_{1,3}
  \begin{vmat}
    0 &1 &0 \\
    0 &0 &1 \\
    1 &0 &0
  \end{vmat}           \\
  &\quad+
  t_{3,1}t_{1,2}t_{2,3}
  \begin{vmat}
    0 &0 &1 \\
    1 &0 &0 \\
    0 &1 &0
  \end{vmat}
  +t_{3,1}t_{2,2}t_{1,3}
  \begin{vmat}
    0 &0 &1 \\
    0 &1 &0 \\
    1 &0 &0
  \end{vmat}
\end{align*}
Pertama-tama bandingkan keenam hasil kali $t$. 
Hasil kali dalam ekspansi~$T$ sama dengan yang terdapat dalam ekspansi 
transposnya; sebagai contoh, $t_{1,2}t_{2,3}t_{3,1}$ berada di bagian atas, sedangkan
$t_{3,1}t_{1,2}t_{2,3}$ berada di bagian bawah.
Hal ini sepenuhnya masuk akal\Dash keenam bentuk di bagian atas berasal dari semua cara 
memilih satu entri~$T$ dari setiap baris dan kolom, sedangkan keenam bentuk di bagian bawah
merupakan semua cara memilih satu entri~$T$ dari setiap kolom dan baris,
sehingga tentu saja keduanya merupakan himpunan yang sama.
%</pf:DetMatrixEqualsDetTrans1>

%<*pf:DetMatrixEqualsDetTrans2>
Selanjutnya perhatikan bahwa dalam kedua ekspansi,
setiap ekspresi hasil kali $t$ belum tentu dipasangkan dengan 
matriks permutasi yang sama.
Sebagai contoh, di bagian atas, $t_{1,2}t_{2,3}t_{3,1}$ dipasangkan dengan
matriks untuk pemetaan 
$1\mapsto 2$, $2\mapsto 3$, $3\mapsto 1$.
Di bagian bawah, $t_{3,1}t_{1,2}t_{2,3}$ 
dipasangkan dengan matriks untuk pemetaan
$1\mapsto 3$, $2\mapsto 1$, $3\mapsto 2$.
Pemetaan kedua merupakan invers dari pemetaan pertama.
Hal ini juga sepenuhnya masuk akal\Dash baik 
transpos matriks maupun invers pemetaan membalik
$1,2$ menjadi $2,1$, membalik $2,3$ menjadi $3,2$, dan membalik $3,1$ menjadi $1,3$.
 
Sebagai penutup, kita mencatat bahwa determinan $P_{\phi}$ sama dengan determinan
$P_{\phi^{-1}}$, sebagaimana ditunjukkan oleh
\nearbyexercise{exer:PermInvFacts}. %
%</pf:DetMatrixEqualsDetTrans2>
\end{proof}


\begin{exercises}
  \item[]\textit{Berikut merangkum notasi yang digunakan dalam buku ini
    untuk permutasi-$2$\hbox{} dan permutasi-$3$.
    \begin{center}
      \begin{tabular}[t]{c|cc}
        $i$          &$1$      &$2$    \\
        \hline
        $\phi_1(i)$  &$1$      &$2$     \\
        $\phi_2(i)$  &$2$      &$1$     
      \end{tabular}
      \qquad
      \begin{tabular}[t]{c|ccc}
        $i$          &$1$     &$2$   &$3$    \\
        \hline
        $\phi_1(i)$  &$1$     &$2$   &$3$    \\
        $\phi_2(i)$  &$1$     &$3$   &$2$    \\
        $\phi_3(i)$  &$2$     &$1$   &$3$    \\
        $\phi_4(i)$  &$2$     &$3$   &$1$    \\
        $\phi_5(i)$  &$3$     &$1$   &$2$    \\
        $\phi_6(i)$  &$3$     &$2$   &$1$    
      \end{tabular}
    \end{center}  }
  \item 
    Berikan ekspansi permutasi suatu matriks umum $\nbyn{2}$
    beserta transposnya.
    \begin{answer}
      Berikut adalah ekspansi permutasi determinan suatu 
      matriks $\nbyn{2}$,
      \begin{equation*}
        \begin{vmat}
          a  &b  \\
          c  &d
        \end{vmat}
         =ad\cdot \begin{vmat}[r]
                    1  &0  \\
                    0  &1
                  \end{vmat}
         +bc\cdot \begin{vmat}[r]
                    0  &1 \\
                    1  &0
                  \end{vmat}
      \end{equation*}
      dan berikut ekspansi permutasi determinan transposnya.
      \begin{equation*}
         \begin{vmat}
            a  &c  \\
            b  &d
         \end{vmat}
         =ad\cdot\begin{vmat}[r]
                    1  &0  \\
                    0  &1
                 \end{vmat}
         +cb\cdot\begin{vmat}[r]
                   0  &1 \\
                   1  &0
                 \end{vmat}
       \end{equation*}
       Seperti pada ekspansi $\nbyn{3}$ yang dijelaskan dalam subbagian ini,
       matriks permutasi dari suku-suku yang bersesuaian merupakan
       transpos satu sama lain (meskipun hal ini tersamar oleh fakta bahwa masing-masing
       sama dengan transposnya sendiri).
    \end{answer}
  \recommended \item 
    \textit{Soal ini juga muncul dalam subbagian sebelumnya.}
    \begin{exparts}
      \partsitem Tentukan invers setiap permutasi-$2$.
      \partsitem Tentukan invers setiap permutasi-$3$.
    \end{exparts}
    \begin{answer}
      Masing-masing hasil berikut mudah diperiksa.
      \begin{exparts}
        \partsitem 
          \begin{tabular}[t]{r|cc}
            \textit{permutasi} &$\phi_1$  &$\phi_2$ \\
             \hline
            \textit{invers}     &$\phi_1$  &$\phi_2$ 
          \end{tabular}
        \partsitem 
          \begin{tabular}[t]{r|cccccc}
            \textit{permutasi} 
              &$\phi_1$ &$\phi_2$ &$\phi_3$ &$\phi_4$ &$\phi_5$ &$\phi_6$ \\
            \hline
            \textit{invers}
              &$\phi_1$ &$\phi_2$ &$\phi_3$ &$\phi_5$ &$\phi_4$ &$\phi_6$ 
          \end{tabular}
      \end{exparts}
    \end{answer}
  \recommended \item
    \begin{exparts}
      \partsitem Tentukan signum setiap permutasi-$2$.
      \partsitem Tentukan signum setiap permutasi-$3$.
    \end{exparts}
    \begin{answer}
      \begin{exparts}
        \partsitem \( \sgn(\phi_1)=+1 \), \( \sgn(\phi_2)=-1 \)
        \partsitem \( \sgn(\phi_1)=+1 \), \( \sgn(\phi_2)=-1 \),
              \( \sgn(\phi_3)=-1 \), \( \sgn(\phi_4)=+1 \),
              \( \sgn(\phi_5)=+1 \), \( \sgn(\phi_6)=-1 \)
      \end{exparts}  
     \end{answer}
  \item 
    Tentukan satu-satunya suku tak nol dalam ekspansi permutasi
    matriks berikut.
    \begin{equation*}
      \begin{vmat}[r]
        0  &1  &0  &0  \\
        1  &0  &1  &0  \\
        0  &1  &0  &1  \\
        0  &0  &1  &0
      \end{vmat}
    \end{equation*}
    Hitung determinan tersebut dengan menentukan signum permutasi yang
    bersesuaian.
    \begin{answer}
      Untuk memperoleh suku tak nol dalam ekspansi permutasi, kita harus menggunakan
      entri \( 1,2 \) dan entri \( 4,3 \).
      Setelah menetapkan kedua entri tersebut, kita juga harus menggunakan entri \( 2,1 \) dan
      entri \( 3,4 \).
      Signum \( \sequence{2,1,4,3} \) adalah \( +1 \) karena dari
      \begin{equation*}
        \begin{mat}[r]
          0  &1  &0  &0  \\
          1  &0  &0  &0  \\
          0  &0  &0  &1  \\
          0  &0  &1  &0            
        \end{mat}
      \end{equation*}
      dua pertukaran baris $\rho_1\leftrightarrow\rho_2$ dan 
      $\rho_3\leftrightarrow\rho_4$ akan menghasilkan matriks identitas.  
    \end{answer}
  \item  
    \cite{Strang}
    Berapakah signum permutasi-$n$
    \( \phi=\sequence{n,n-1,\dots,2,1} \)?
    \begin{answer}
      Polanya adalah sebagai berikut.
      \begin{center}
         \begin{tabular}[b]{c|ccccccc}
           \( i \) &\( 1 \) &\( 2 \) &\( 3 \) &\( 4 \) &\( 5 \)
              &\( 6 \) &\ldots \\
           \hline
           \( \sgn(\phi_i) \) &\( +1 \) &\( -1 \) &\( -1 \) &\( +1 \)
              &\( +1 \) &\( -1 \) &\ldots
         \end{tabular}
      \end{center}  
      Permutasi tersebut memiliki
      $1+2+\cdots+(n-1)=n(n-1)/2$ inversi.
      Karena itu, signumnya adalah
      $(-1)^{n(n-1)/2}$.
      Dengan kata lain, signumnya $+1$ jika $n$ bersisa $0$ atau $1$ ketika
      dibagi empat, dan $-1$ jika $n$ bersisa $2$ atau $3$ ketika dibagi empat.
    \end{answer}
  \item \label{exer:PermInvFacts}
     Buktikan pernyataan-pernyataan berikut.   
     \begin{exparts}
        \partsitem Setiap permutasi memiliki invers.
        \partsitem \( \sgn(\phi^{-1})=\sgn(\phi) \)
        \partsitem Setiap permutasi merupakan invers dari suatu permutasi lain.
     \end{exparts}
    \begin{answer}
      \begin{exparts}
       \partsitem Kita dapat memandang permutasi sebagai pemetaan
         \( \map{\phi}{\set{1,\ldots,n}}{\set{1,\ldots,n}} \)
         yang bersifat satu-ke-satu dan pada.
         Setiap pemetaan satu-ke-satu dan pada memiliki invers.
       \partsitem Jika selalu diperlukan banyak pertukaran ganjil untuk mengubah 
         \( P_\phi \)
         menjadi matriks identitas, maka selalu diperlukan banyak pertukaran ganjil
         untuk mengubah matriks identitas menjadi \( P_\phi \) (setiap pertukaran dapat dibalik).
       \partsitem Ini merupakan pertanyaan pertama lagi.
     \end{exparts}  
    \end{answer}
  \item \label{exer:PermInvIffMatTrans}
      Buktikan bahwa matriks dari invers suatu permutasi
      merupakan transpos matriks permutasi tersebut, 
      $P_{\phi^{-1}}=\trans{P_{\phi}}$, 
      untuk setiap permutasi $\phi$.
      \begin{answer}
        Jika $\phi(i)=j$, maka $\phi^{-1}(j)=i$.
        Hasilnya kini diperoleh dengan mengamati bahwa $P_{\phi}$ memiliki 
        $1$ pada entri $i,j$ jika dan hanya jika $\phi(i)=j$,
        dan $P_{\phi^{-1}}$ memiliki 
        $1$ pada entri $j,i$ jika dan hanya jika $\phi^{-1}(j)=i$.
      \end{answer}
  \recommended \item 
    Tunjukkan bahwa matriks permutasi dengan \( m \) inversi dapat
    diubah melalui pertukaran baris menjadi matriks identitas dalam \( m \) langkah.
    Bandingkan hal ini dengan \nearbycorollary{cor:ParityInversEqParitySwaps}.
    \begin{answer}
      Pernyataan ini tidak mengatakan bahwa \( m \) merupakan jumlah pertukaran terkecil untuk menghasilkan
      matriks identitas, dan juga tidak mengatakan bahwa \( m \) merupakan jumlah terbesar.
      Pernyataan tersebut justru mengatakan bahwa 
      terdapat cara untuk melakukan pertukaran hingga menjadi matriks identitas dalam tepat \( m \) langkah.

      Misalkan \( \iota_j \) adalah baris pertama yang mengalami inversi relatif
      terhadap suatu baris sebelumnya,
      dan misalkan \( \iota_k \) adalah baris pertama yang menimbulkan inversi tersebut.
      Kita memiliki interval baris berikut.
      \begin{equation*}
         \begin{mat}
           \vdots      \\
           \iota_k     \\
           \iota_{r_1} \\
           \vdots      \\
           \iota_{r_s} \\
           \iota_j     \\
           \vdots
         \end{mat}
         \qquad j<k<r_1<\dots <r_s
      \end{equation*}
      Lakukan pertukaran.
      \begin{equation*}
         \begin{mat}
           \vdots      \\
           \iota_j     \\
           \iota_{r_1} \\
           \vdots      \\
           \iota_{r_s} \\
           \iota_k     \\
           \vdots
         \end{mat}
      \end{equation*}
      Matriks kedua memiliki satu inversi lebih sedikit karena terdapat satu
      inversi lebih sedikit
      dalam interval tersebut (\( s \) dibandingkan dengan\ \( s+1 \)), sedangkan inversi yang melibatkan
      baris di luar interval tidak terpengaruh.

      Lanjutkan dengan cara ini; pada setiap langkah, kurangi banyaknya inversi
      sebanyak satu
      melalui setiap pertukaran baris.
      Ketika tidak ada inversi yang tersisa, hasilnya adalah matriks identitas. 

      Perbedaannya dengan \nearbycorollary{cor:ParityInversEqParitySwaps}
      ialah bahwa pernyataan soal ini merupakan pernyataan `terdapat':
      terdapat cara untuk melakukan pertukaran hingga menjadi matriks identitas dalam tepat~$m$ langkah.
      Namun, korolari tersebut merupakan pernyataan `untuk semua': untuk semua cara melakukan pertukaran hingga menjadi
      matriks identitas, paritasnya (genap atau ganjil) sama. 
    \end{answer}
  \recommended \item 
    Untuk setiap permutasi $\phi$, misalkan $g(\phi)$ merupakan bilangan bulat 
    yang didefinisikan sebagai berikut.
    \begin{equation*}
       g(\phi)=\prod_{i<j} [\phi(j)-\phi(i)]
    \end{equation*}
    (Ini merupakan hasil kali, atas semua indeks \( i \) dan \( j \) dengan \( i<j \),
    dari suku-suku berbentuk seperti yang diberikan.)
    \begin{exparts}
      \partsitem Hitung nilai $g$ pada semua permutasi-$2$.
      \partsitem Hitung nilai $g$ pada semua permutasi-$3$.
      \partsitem Buktikan bahwa $g(\phi)$ tidak bernilai $0$.
      \partsitem Buktikan pernyataan berikut.
         \begin{equation*}
            \sgn(\phi)=\frac{g(\phi)}{\absval{g(\phi)}}
         \end{equation*}
    \end{exparts}
    Banyak penulis memberikan rumus ini sebagai definisi fungsi signum.
    \begin{answer}
      \begin{exparts}
        \partsitem Pertama, $g(\phi_1)$ merupakan hasil kali faktor tunggal
          $2-1$, sehingga $g(\phi_1)=1$.
          Kedua, $g(\phi_2)$ merupakan hasil kali faktor tunggal
          $1-2$, sehingga $g(\phi_2)=-1$.
        \partsitem 
            \begin{tabular}[t]{r|cccccc}
               \textit{permutasi $\phi$} 
                 &$\phi_{1}$ &$\phi_{2}$ &$\phi_{3}$ 
                    &$\phi_{4}$ &$\phi_{5}$ &$\phi_{6}$ \\
               \hline
               $g(\phi)$ &$2$ &$-2$ &$-2$ &$2$ &$2$ &$-2$ 
            \end{tabular}
        \partsitem Bentuk tersebut merupakan hasil kali suku-suku tak nol.
        \partsitem Perhatikan bahwa \( \phi(j)-\phi(i) \) bernilai negatif jika dan hanya jika
           \( \iota_{\phi(j)} \) dan \( \iota_{\phi(i)} \) berada dalam suatu inversi
           dari urutan biasanya.
      \end{exparts}  
     \end{answer}
%  \item Prove that the signum function on permutations is unique in that
%    if \( f \) is a map from the permutations of \( \set{1,\ldots,n} \)
%    to the integers that is multiplicative:
%    \( f(\sigma\phi)=f(\sigma)f(\phi) \) then \( f \) is either the
%    function that always gives zero, or the function that always gives
%    one, or is the signum.
%    \begin{answer}
%      Consider the identity permutation.
%      \begin{equation*}
%        \map{\identity}{\set{1,2,\ldots,n}}{\set{1,2,\ldots,n}}    
%      \end{equation*}
%      Note that $f(\identity)$ equals $0$ or $1$ because
%      the multiplicative property gives 
%      $f(\identity)
%       =f(\composed{\identity}{\identity})
%       =f(\identity)\cdot f(\identity)
%       =f(\identity)^2$.
%      If $f(\identity)=0$ then $f$ is the constant function 
%      $f(\phi)=0$ because 
%      $f(\phi)=f(\composed{\identity}{\phi})=0\cdot f(\phi)=0$.
%
%      If $f(\identity)=1$ then consider the permutation
%      $\tau$ that just swaps $i$ and $j$
%      \begin{equation*}
%        1\mapsto 1
%        \quad 2\mapsto 2
%        \quad \cdots 
%        \quad i\mapsto j
%        \quad \cdots
%        \quad j\mapsto i
%        \quad \cdots
%        \quad n\mapsto n
%      \end{equation*}
%      (the $n=1$ special case is easy).
%      Note that $\composed{\tau}{\tau}=\identity$ so 
%      $1=f(\composed{\tau}{\tau})=f(\tau)\cdot f(\tau)$,
%      and $f(\tau)$ is either $+1$ or $-1$.
%    \end{answer}
\end{exercises}



\index{determinan|)}
